
\documentclass{article}
\usepackage{colortbl}
\usepackage{makecell}
\usepackage{multirow}
\usepackage{supertabular}

\begin{document}

\newcounter{utterance}

\twocolumn

{ \footnotesize  \setcounter{utterance}{1}
\setlength{\tabcolsep}{0pt}
\begin{supertabular}{c@{$\;$}|p{.15\linewidth}@{}p{.15\linewidth}p{.15\linewidth}p{.15\linewidth}p{.15\linewidth}p{.15\linewidth}}

    \# & $\;$A & \multicolumn{4}{c}{Game Master} & $\;\:$B\\
    \hline 

    \theutterance \stepcounter{utterance}  

    & & \multicolumn{4}{p{0.6\linewidth}}{\cellcolor[rgb]{0.9,0.9,0.9}{%
	\makecell[{{p{\linewidth}}}]{% 
	  \tt {\tiny [A$\langle$GM]}  
	 Sono il tuo game master, e stai giocando un gioco collaborativo con la seguente griglia come tavolo da gioco:\\ \tt \\ \tt ```\\ \tt  1234567\\ \tt ╔═╤═╤═╤═╗\\ \tt ║◌│◌│H│◌║ 1\\ \tt ╟─┴─┼─┤◌║ 2\\ \tt ║◌T◌│◌└─╢ 3\\ \tt ╟───┘◌◌◌║ 4\\ \tt ║◌◌◌I◌┌─╢ 5\\ \tt ║◌◌◌C◌│◌║ 6\\ \tt ║◌◌◌W◌│◌║ 7\\ \tt ╚═════╧═╝\\ \tt \\ \tt ```\\ \tt \\ \tt * Il bordo superiore mostra le coordinate x che aumentano verso destra, e il bordo destro le coordinate y che aumentano verso il basso.\\ \tt * The following objects are randomly placed on your grid: 'W', 'I', 'T', 'C', 'H'.\\ \tt \\ \tt \\ \tt L'altro giocatore vede una variante del tavolo da gioco, dove gli oggetti sono posizionati in diverse posizioni casuali. Non puoi vedere il tavolo dell'altro giocatore, e loro non possono vedere il tuo.\\ \tt \\ \tt \\ \tt **Obiettivo:**\\ \tt \\ \tt Entrambi i giocatori devono spostare gli oggetti sui rispettivi sfondi in modo che oggetti identici finiscano alle stesse coordinate. Devi comunicare con l'altro giocatore per concordare una configurazione obiettivo comune.\\ \tt \\ \tt \\ \tt **Regole:**\\ \tt \\ \tt * In ogni turno, puoi inviare esattamente uno dei seguenti due comandi:\\ \tt  1. `COMUNICA: <MESSAGE>`: per inviare un messaggio (tutto fino alla prossima interruzione di riga) all'altro giocatore. Lo inoltrerò al tuo partner.\\ \tt 2. `SPOSTA: <OBJECT>, (<X>, <Y>)`: per spostare un oggetto in una nuova posizione, dove <X> è la colonna e <Y> è la riga. Ti farò sapere se la tua mossa è andata a buon fine o no.\\ \tt * Se non rispetti il formato, o invii diversi comandi contemporaneamente, devo penalizzarti.\\ \tt * If both players accumulate more than 12 penalties, you both lose the game.\\ \tt * È vitale che comunichi con l'altro giocatore riguardo al tuo stato obiettivo! L'*unico* modo in cui puoi trasmettere la tua strategia all'altro giocatore è usando il comando `COMUNICA: <MESSAGE>`!\\ \tt \\ \tt \\ \tt **Spostamento Oggetti**\\ \tt \\ \tt * Puoi spostare oggetti solo in celle entro i limiti della griglia. La cella di destinazione deve essere vuota, cioè deve contenere solo il simbolo '◌'.\\ \tt * Se provi a spostare un oggetto in un posto che non è vuoto, o provi a spostarlo fuori dalla griglia, devo penalizzarti. Hai un altro tentativo.\\ \tt * Prima di fare una mossa, controlla due volte che il posto di destinazione sia vuoto, e non contenga alcuna lettera, cornice, o linea!\\ \tt \\ \tt \\ \tt **Fine del Gioco**\\ \tt \\ \tt Se pensi di aver raggiunto l'obiettivo di allineare tutti gli oggetti, puoi chiedere all'altro giocatore di finire il gioco inviando COMUNICA: finito?. Se l'altro giocatore ti chiede di finire il gioco, e rispondi con COMUNICA: finito!, the game will end.\\ \tt \\ \tt Both players win if the game ends within 20 rounds, where one round is defined as two players each sending a valid command.\\ \tt \\ \tt **Punteggio:**\\ \tt \\ \tt Più vicini sono gli oggetti identici in entrambi i tavoli da gioco, più punti ottieni. Le penalità riducono i tuoi punti. Riesci a battere il record?\\ \tt \\ \tt \\ \tt Per favore invia un messaggio all'altro giocatore per iniziare il gioco! 
	  } 
	   } 
	   } 
	 & \\ 
 

    \theutterance \stepcounter{utterance}  

    & \multicolumn{4}{p{0.6\linewidth}}{\cellcolor[rgb]{0.8,1,0.9}{%
	\makecell[{{p{\linewidth}}}]{% 
	  \tt {\tiny [A$\rangle$GM]}  
	 Sembra un gioco interessante! Iniziamo comunicando la nostra strategia per allineare gli oggetti. Ecco il mio messaggio di inizio:\\ \tt \\ \tt COMUNICA: Ciao! Per iniziare, vorrei proporre di sistemare gli oggetti in base all'alfabeto, partendo dalla posizione più in alto a sinistra. Come pensi di procedere? Qual è la tua idea per allineare gli oggetti? 
	  } 
	   } 
	   } 
	 & & \\ 
 

    \theutterance \stepcounter{utterance}  

    & & & \multicolumn{2}{p{0.3\linewidth}}{\cellcolor[rgb]{0.95,0.95,0.95}{%
	\makecell[{{p{\linewidth}}}]{% 
	  \tt {\tiny [GM$|$GM]}  
	 Sembra un gioco interessante! Iniziamo comunicando la nostra strategia per allineare gli oggetti. Ecco il mio messaggio di inizio:\\ \tt \\ \tt COMUNICA: Ciao! Per iniziare, vorrei proporre di sistemare gli oggetti in base all'alfabeto, partendo dalla posizione più in alto a sinistra. Come pensi di procedere? Qual è la tua idea per allineare gli oggetti? 
	  } 
	   } 
	   } 
	 & & \\ 
 

    \theutterance \stepcounter{utterance}  

    & & & \multicolumn{2}{p{0.3\linewidth}}{\cellcolor[rgb]{0.95,0.95,0.95}{%
	\makecell[{{p{\linewidth}}}]{% 
	  \tt {\tiny [GM$|$GM]}  
	 Invalid format: head is not empty: 'Sembra un gioco interessante! Iniziamo comunicando la nostra strategia per allineare gli oggetti. Ecco il mio messaggio di inizio:\\ \tt \\ \tt ' 
	  } 
	   } 
	   } 
	 & & \\ 
 

    \theutterance \stepcounter{utterance}  

    & & \multicolumn{4}{p{0.6\linewidth}}{\cellcolor[rgb]{0.9,0.9,0.9}{%
	\makecell[{{p{\linewidth}}}]{% 
	  \tt {\tiny [A$\langle$GM]}  
	 Penalità: Il tuo messaggio non deve contenere nulla prima del comando!\\ \tt \\ \tt Assicurati che la tua risposta contenga solo `COMUNICA: <MESSAGE>` o `SPOSTA: <OBJECT>, (<X>, <Y>)`, e nient'altro!\\ \tt \\ \tt You have collectively accumulated 1 of 12 penalties.\\ \tt Riprova! 
	  } 
	   } 
	   } 
	 & \\ 
 

    \theutterance \stepcounter{utterance}  

    & \multicolumn{4}{p{0.6\linewidth}}{\cellcolor[rgb]{0.8,1,0.9}{%
	\makecell[{{p{\linewidth}}}]{% 
	  \tt {\tiny [A$\rangle$GM]}  
	 COMUNICA: Ciao! Per iniziare, vorrei proporre di sistemare gli oggetti in base all'alfabeto, partendo dalla posizione più in alto a sinistra. Come pensi di procedere? Qual è la tua idea per allineare gli oggetti? 
	  } 
	   } 
	   } 
	 & & \\ 
 

    \theutterance \stepcounter{utterance}  

    & & & \multicolumn{2}{p{0.3\linewidth}}{\cellcolor[rgb]{0.95,0.95,0.95}{%
	\makecell[{{p{\linewidth}}}]{% 
	  \tt {\tiny [GM$|$GM]}  
	 COMUNICA: Ciao! Per iniziare, vorrei proporre di sistemare gli oggetti in base all'alfabeto, partendo dalla posizione più in alto a sinistra. Come pensi di procedere? Qual è la tua idea per allineare gli oggetti? 
	  } 
	   } 
	   } 
	 & & \\ 
 

    \theutterance \stepcounter{utterance}  

    & & \multicolumn{4}{p{0.6\linewidth}}{\cellcolor[rgb]{0.9,0.9,0.9}{%
	\makecell[{{p{\linewidth}}}]{% 
	  \tt {\tiny [GM$\rangle$B]}  
	 Sono il tuo game master, e stai giocando un gioco collaborativo con la seguente griglia come tavolo da gioco:\\ \tt \\ \tt ```\\ \tt  1234567\\ \tt ╔═╤═╤═╤═╗\\ \tt ║W│◌│T│◌║ 1\\ \tt ╟─┴─┼─┤◌║ 2\\ \tt ║◌◌◌│◌└─╢ 3\\ \tt ╟───┘IC◌║ 4\\ \tt ║H◌◌◌◌┌─╢ 5\\ \tt ║◌◌◌◌◌│◌║ 6\\ \tt ║◌◌◌◌◌│◌║ 7\\ \tt ╚═════╧═╝\\ \tt \\ \tt ```\\ \tt \\ \tt * Il bordo superiore mostra le coordinate x che aumentano verso destra, e il bordo destro le coordinate y che aumentano verso il basso.\\ \tt * The following objects are randomly placed on your grid: 'W', 'I', 'T', 'C', 'H'.\\ \tt \\ \tt \\ \tt L'altro giocatore vede una variante del tavolo da gioco, dove gli oggetti sono posizionati in diverse posizioni casuali. Non puoi vedere il tavolo dell'altro giocatore, e loro non possono vedere il tuo.\\ \tt \\ \tt \\ \tt **Obiettivo:**\\ \tt \\ \tt Entrambi i giocatori devono spostare gli oggetti sui rispettivi sfondi in modo che oggetti identici finiscano alle stesse coordinate. Devi comunicare con l'altro giocatore per concordare una configurazione obiettivo comune.\\ \tt \\ \tt \\ \tt **Regole:**\\ \tt \\ \tt * In ogni turno, puoi inviare esattamente uno dei seguenti due comandi:\\ \tt  1. `COMUNICA: <MESSAGE>`: per inviare un messaggio (tutto fino alla prossima interruzione di riga) all'altro giocatore. Lo inoltrerò al tuo partner.\\ \tt 2. `SPOSTA: <OBJECT>, (<X>, <Y>)`: per spostare un oggetto in una nuova posizione, dove <X> è la colonna e <Y> è la riga. Ti farò sapere se la tua mossa è andata a buon fine o no.\\ \tt * Se non rispetti il formato, o invii diversi comandi contemporaneamente, devo penalizzarti.\\ \tt * If both players accumulate more than 12 penalties, you both lose the game.\\ \tt * È vitale che comunichi con l'altro giocatore riguardo al tuo stato obiettivo! L'*unico* modo in cui puoi trasmettere la tua strategia all'altro giocatore è usando il comando `COMUNICA: <MESSAGE>`!\\ \tt \\ \tt \\ \tt **Spostamento Oggetti**\\ \tt \\ \tt * Puoi spostare oggetti solo in celle entro i limiti della griglia. La cella di destinazione deve essere vuota, cioè deve contenere solo il simbolo '◌'.\\ \tt * Se provi a spostare un oggetto in un posto che non è vuoto, o provi a spostarlo fuori dalla griglia, devo penalizzarti. Hai un altro tentativo.\\ \tt * Prima di fare una mossa, controlla due volte che il posto di destinazione sia vuoto, e non contenga alcuna lettera, cornice, o linea!\\ \tt \\ \tt \\ \tt **Fine del Gioco**\\ \tt \\ \tt Se pensi di aver raggiunto l'obiettivo di allineare tutti gli oggetti, puoi chiedere all'altro giocatore di finire il gioco inviando COMUNICA: finito?. Se l'altro giocatore ti chiede di finire il gioco, e rispondi con COMUNICA: finito!, the game will end.\\ \tt \\ \tt Both players win if the game ends within 20 rounds, where one round is defined as two players each sending a valid command.\\ \tt \\ \tt **Punteggio:**\\ \tt \\ \tt Più vicini sono gli oggetti identici in entrambi i tavoli da gioco, più punti ottieni. Le penalità riducono i tuoi punti. Riesci a battere il record?\\ \tt \\ \tt L'altro giocatore ha iniziato il gioco inviando questo messaggio:\\ \tt \\ \tt "Ciao! Per iniziare, vorrei proporre di sistemare gli oggetti in base all'alfabeto, partendo dalla posizione più in alto a sinistra. Come pensi di procedere? Qual è la tua idea per allineare gli oggetti?"\\ \tt \\ \tt Qual è il tuo primo comando? 
	  } 
	   } 
	   } 
	 & \\ 
 

    \theutterance \stepcounter{utterance}  

    & & & \multicolumn{4}{p{0.6\linewidth}}{\cellcolor[rgb]{1,0.85,0.72}{%
	\makecell[{{p{\linewidth}}}]{% 
	  \tt {\tiny [GM$\langle$B]}  
	 COMUNICA: Ciao! La tua idea di sistemare gli oggetti in base all'alfabeto mi sembra interessante. Per procedere, potremmo iniziare a scambiare informazioni sulla posizione degli oggetti sul nostro tavolo. Io ho un 'W' in alto a sinistra, un 'T' nella seconda colonna della prima riga, un 'I' e un 'C' e un 'H' in posizioni diverse. Potresti darmi un'idea di come sono posizionati i tuoi oggetti? Inoltre, come pensi di gestire gli spostamenti degli oggetti per allinearli? Dovremmo iniziare con gli oggetti che sono già vicini alle loro posizioni di destinazione? 
	  } 
	   } 
	   } 
	  \\ 
 

    \theutterance \stepcounter{utterance}  

    & & & \multicolumn{2}{p{0.3\linewidth}}{\cellcolor[rgb]{0.95,0.95,0.95}{%
	\makecell[{{p{\linewidth}}}]{% 
	  \tt {\tiny [GM$|$GM]}  
	 COMUNICA: Ciao! La tua idea di sistemare gli oggetti in base all'alfabeto mi sembra interessante. Per procedere, potremmo iniziare a scambiare informazioni sulla posizione degli oggetti sul nostro tavolo. Io ho un 'W' in alto a sinistra, un 'T' nella seconda colonna della prima riga, un 'I' e un 'C' e un 'H' in posizioni diverse. Potresti darmi un'idea di come sono posizionati i tuoi oggetti? Inoltre, come pensi di gestire gli spostamenti degli oggetti per allinearli? Dovremmo iniziare con gli oggetti che sono già vicini alle loro posizioni di destinazione? 
	  } 
	   } 
	   } 
	 & & \\ 
 

    \theutterance \stepcounter{utterance}  

    & & \multicolumn{4}{p{0.6\linewidth}}{\cellcolor[rgb]{0.9,0.9,0.9}{%
	\makecell[{{p{\linewidth}}}]{% 
	  \tt {\tiny [A$\langle$GM]}  
	 Il tuo messaggio è stato inoltrato all'altro giocatore.\\ \tt \\ \tt \\ \tt You are currently playing round 1 of maximum 20.\\ \tt You have collectively accumulated 1 of 12 penalties.\\ \tt L'altro giocatore ha inviato questo messaggio:\\ \tt \\ \tt \\ \tt "Ciao! La tua idea di sistemare gli oggetti in base all'alfabeto mi sembra interessante. Per procedere, potremmo iniziare a scambiare informazioni sulla posizione degli oggetti sul nostro tavolo. Io ho un 'W' in alto a sinistra, un 'T' nella seconda colonna della prima riga, un 'I' e un 'C' e un 'H' in posizioni diverse. Potresti darmi un'idea di come sono posizionati i tuoi oggetti? Inoltre, come pensi di gestire gli spostamenti degli oggetti per allinearli? Dovremmo iniziare con gli oggetti che sono già vicini alle loro posizioni di destinazione?"\\ \tt \\ \tt \\ \tt Qual è il tuo prossimo comando? 
	  } 
	   } 
	   } 
	 & \\ 
 

    \theutterance \stepcounter{utterance}  

    & \multicolumn{4}{p{0.6\linewidth}}{\cellcolor[rgb]{0.8,1,0.9}{%
	\makecell[{{p{\linewidth}}}]{% 
	  \tt {\tiny [A$\rangle$GM]}  
	 COMUNICA: Grazie per le informazioni! Io ho un 'H' nella prima riga, un 'T' nella terza riga, un 'I' e un 'C' e un 'W' in posizioni diverse. Penso che dovremmo iniziare con gli oggetti che sono già vicini alle loro posizioni di destinazione, come hai suggerito. Per gestire gli spostamenti, potremmo iniziare a spostare gli oggetti uno alla volta, partendo da quelli che sono più vicini alla loro posizione di destinazione. Potremmo anche utilizzare una convenzione per indicare le posizioni degli oggetti, ad esempio utilizzando le coordinate (x, y) per indicare la posizione di ogni oggetto. Cosa ne pensi? 
	  } 
	   } 
	   } 
	 & & \\ 
 

    \theutterance \stepcounter{utterance}  

    & & & \multicolumn{2}{p{0.3\linewidth}}{\cellcolor[rgb]{0.95,0.95,0.95}{%
	\makecell[{{p{\linewidth}}}]{% 
	  \tt {\tiny [GM$|$GM]}  
	 COMUNICA: Grazie per le informazioni! Io ho un 'H' nella prima riga, un 'T' nella terza riga, un 'I' e un 'C' e un 'W' in posizioni diverse. Penso che dovremmo iniziare con gli oggetti che sono già vicini alle loro posizioni di destinazione, come hai suggerito. Per gestire gli spostamenti, potremmo iniziare a spostare gli oggetti uno alla volta, partendo da quelli che sono più vicini alla loro posizione di destinazione. Potremmo anche utilizzare una convenzione per indicare le posizioni degli oggetti, ad esempio utilizzando le coordinate (x, y) per indicare la posizione di ogni oggetto. Cosa ne pensi? 
	  } 
	   } 
	   } 
	 & & \\ 
 

    \theutterance \stepcounter{utterance}  

    & & \multicolumn{4}{p{0.6\linewidth}}{\cellcolor[rgb]{0.9,0.9,0.9}{%
	\makecell[{{p{\linewidth}}}]{% 
	  \tt {\tiny [GM$\rangle$B]}  
	 Il tuo messaggio è stato inoltrato all'altro giocatore.\\ \tt \\ \tt \\ \tt You are currently playing round 2 of maximum 20.\\ \tt You have collectively accumulated 1 of 12 penalties.\\ \tt L'altro giocatore ha inviato questo messaggio:\\ \tt \\ \tt \\ \tt "Grazie per le informazioni! Io ho un 'H' nella prima riga, un 'T' nella terza riga, un 'I' e un 'C' e un 'W' in posizioni diverse. Penso che dovremmo iniziare con gli oggetti che sono già vicini alle loro posizioni di destinazione, come hai suggerito. Per gestire gli spostamenti, potremmo iniziare a spostare gli oggetti uno alla volta, partendo da quelli che sono più vicini alla loro posizione di destinazione. Potremmo anche utilizzare una convenzione per indicare le posizioni degli oggetti, ad esempio utilizzando le coordinate (x, y) per indicare la posizione di ogni oggetto. Cosa ne pensi?"\\ \tt \\ \tt \\ \tt Qual è il tuo prossimo comando? 
	  } 
	   } 
	   } 
	 & \\ 
 

    \theutterance \stepcounter{utterance}  

    & & & \multicolumn{4}{p{0.6\linewidth}}{\cellcolor[rgb]{1,0.85,0.72}{%
	\makecell[{{p{\linewidth}}}]{% 
	  \tt {\tiny [GM$\langle$B]}  
	 COMUNICA: Sono d'accordo con la tua proposta di iniziare con gli oggetti vicini alle loro posizioni di destinazione e di utilizzare le coordinate (x, y) per indicare le posizioni. Per facilitare la comunicazione, potremmo anche stabilire un ordine standard per elencare gli oggetti, ad esempio in base all'alfabeto ('C', 'H', 'I', 'T', 'W'). \\ \tt \\ \tt Vorrei anche chiederti di confermare se le tue posizioni di 'H' e 'T' sono rispettivamente nella prima e terza riga, come hai detto, e se il tuo 'W' si trova anche in una posizione relativamente alta. Queste informazioni mi aiuteranno a comprendere meglio la disposizione degli oggetti sul tuo tavolo e a pianificare i prossimi spostamenti in modo efficiente. 
	  } 
	   } 
	   } 
	  \\ 
 

    \theutterance \stepcounter{utterance}  

    & & & \multicolumn{2}{p{0.3\linewidth}}{\cellcolor[rgb]{0.95,0.95,0.95}{%
	\makecell[{{p{\linewidth}}}]{% 
	  \tt {\tiny [GM$|$GM]}  
	 COMUNICA: Sono d'accordo con la tua proposta di iniziare con gli oggetti vicini alle loro posizioni di destinazione e di utilizzare le coordinate (x, y) per indicare le posizioni. Per facilitare la comunicazione, potremmo anche stabilire un ordine standard per elencare gli oggetti, ad esempio in base all'alfabeto ('C', 'H', 'I', 'T', 'W'). \\ \tt \\ \tt Vorrei anche chiederti di confermare se le tue posizioni di 'H' e 'T' sono rispettivamente nella prima e terza riga, come hai detto, e se il tuo 'W' si trova anche in una posizione relativamente alta. Queste informazioni mi aiuteranno a comprendere meglio la disposizione degli oggetti sul tuo tavolo e a pianificare i prossimi spostamenti in modo efficiente. 
	  } 
	   } 
	   } 
	 & & \\ 
 

    \theutterance \stepcounter{utterance}  

    & & & \multicolumn{2}{p{0.3\linewidth}}{\cellcolor[rgb]{0.95,0.95,0.95}{%
	\makecell[{{p{\linewidth}}}]{% 
	  \tt {\tiny [GM$|$GM]}  
	 Invalid format: tail is not empty: '\\ \tt \\ \tt Vorrei anche chiederti di confermare se le tue posizioni di 'H' e 'T' sono rispettivamente nella prima e terza riga, come hai detto, e se il tuo 'W' si trova anche in una posizione relativamente alta. Queste informazioni mi aiuteranno a comprendere meglio la disposizione degli oggetti sul tuo tavolo e a pianificare i prossimi spostamenti in modo efficiente. 
	  } 
	   } 
	   } 
	 & & \\ 
 

    \theutterance \stepcounter{utterance}  

    & & \multicolumn{4}{p{0.6\linewidth}}{\cellcolor[rgb]{0.9,0.9,0.9}{%
	\makecell[{{p{\linewidth}}}]{% 
	  \tt {\tiny [GM$\rangle$B]}  
	 Penalità: Il tuo messaggio non deve contenere nulla dopo il comando!\\ \tt \\ \tt Assicurati che la tua risposta contenga solo `COMUNICA: <MESSAGE>` o `SPOSTA: <OBJECT>, (<X>, <Y>)`, e nient'altro!\\ \tt \\ \tt You have collectively accumulated 2 of 12 penalties.\\ \tt Riprova! 
	  } 
	   } 
	   } 
	 & \\ 
 

    \theutterance \stepcounter{utterance}  

    & & & \multicolumn{4}{p{0.6\linewidth}}{\cellcolor[rgb]{1,0.85,0.72}{%
	\makecell[{{p{\linewidth}}}]{% 
	  \tt {\tiny [GM$\langle$B]}  
	 COMUNICA: Sono d'accordo con la tua proposta di iniziare con gli oggetti vicini alle loro posizioni di destinazione e di utilizzare le coordinate (x, y) per indicare le posizioni. Potremmo anche stabilire un ordine standard per elencare gli oggetti, ad esempio in base all'alfabeto ('C', 'H', 'I', 'T', 'W'). Vorrei chiederti di confermare se le tue posizioni di 'H' e 'T' sono rispettivamente nella prima e terza riga, come hai detto, e se il tuo 'W' si trova anche in una posizione relativamente alta. 
	  } 
	   } 
	   } 
	  \\ 
 

    \theutterance \stepcounter{utterance}  

    & & & \multicolumn{2}{p{0.3\linewidth}}{\cellcolor[rgb]{0.95,0.95,0.95}{%
	\makecell[{{p{\linewidth}}}]{% 
	  \tt {\tiny [GM$|$GM]}  
	 COMUNICA: Sono d'accordo con la tua proposta di iniziare con gli oggetti vicini alle loro posizioni di destinazione e di utilizzare le coordinate (x, y) per indicare le posizioni. Potremmo anche stabilire un ordine standard per elencare gli oggetti, ad esempio in base all'alfabeto ('C', 'H', 'I', 'T', 'W'). Vorrei chiederti di confermare se le tue posizioni di 'H' e 'T' sono rispettivamente nella prima e terza riga, come hai detto, e se il tuo 'W' si trova anche in una posizione relativamente alta. 
	  } 
	   } 
	   } 
	 & & \\ 
 

    \theutterance \stepcounter{utterance}  

    & & \multicolumn{4}{p{0.6\linewidth}}{\cellcolor[rgb]{0.9,0.9,0.9}{%
	\makecell[{{p{\linewidth}}}]{% 
	  \tt {\tiny [A$\langle$GM]}  
	 Il tuo messaggio è stato inoltrato all'altro giocatore.\\ \tt \\ \tt \\ \tt You are currently playing round 2 of maximum 20.\\ \tt You have collectively accumulated 2 of 12 penalties.\\ \tt L'altro giocatore ha inviato questo messaggio:\\ \tt \\ \tt \\ \tt "Sono d'accordo con la tua proposta di iniziare con gli oggetti vicini alle loro posizioni di destinazione e di utilizzare le coordinate (x, y) per indicare le posizioni. Potremmo anche stabilire un ordine standard per elencare gli oggetti, ad esempio in base all'alfabeto ('C', 'H', 'I', 'T', 'W'). Vorrei chiederti di confermare se le tue posizioni di 'H' e 'T' sono rispettivamente nella prima e terza riga, come hai detto, e se il tuo 'W' si trova anche in una posizione relativamente alta."\\ \tt \\ \tt \\ \tt Qual è il tuo prossimo comando? 
	  } 
	   } 
	   } 
	 & \\ 
 

    \theutterance \stepcounter{utterance}  

    & \multicolumn{4}{p{0.6\linewidth}}{\cellcolor[rgb]{0.8,1,0.9}{%
	\makecell[{{p{\linewidth}}}]{% 
	  \tt {\tiny [A$\rangle$GM]}  
	 COMUNICA: Sì, confermo che il mio 'H' si trova nella prima riga e il mio 'T' si trova nella terza riga. Il mio 'W' si trova nella settima riga. Sono d'accordo con l'ordine standard di elencare gli oggetti in base all'alfabeto ('C', 'H', 'I', 'T', 'W'). Per procedere, potremmo iniziare a scambiare informazioni sulla posizione esatta di ogni oggetto, utilizzando le coordinate (x, y) per indicare la posizione di ogni oggetto. Ad esempio, potrei dire che il mio 'H' si trova nella posizione (3, 1). Come pensi di procedere? 
	  } 
	   } 
	   } 
	 & & \\ 
 

    \theutterance \stepcounter{utterance}  

    & & & \multicolumn{2}{p{0.3\linewidth}}{\cellcolor[rgb]{0.95,0.95,0.95}{%
	\makecell[{{p{\linewidth}}}]{% 
	  \tt {\tiny [GM$|$GM]}  
	 COMUNICA: Sì, confermo che il mio 'H' si trova nella prima riga e il mio 'T' si trova nella terza riga. Il mio 'W' si trova nella settima riga. Sono d'accordo con l'ordine standard di elencare gli oggetti in base all'alfabeto ('C', 'H', 'I', 'T', 'W'). Per procedere, potremmo iniziare a scambiare informazioni sulla posizione esatta di ogni oggetto, utilizzando le coordinate (x, y) per indicare la posizione di ogni oggetto. Ad esempio, potrei dire che il mio 'H' si trova nella posizione (3, 1). Come pensi di procedere? 
	  } 
	   } 
	   } 
	 & & \\ 
 

    \theutterance \stepcounter{utterance}  

    & & \multicolumn{4}{p{0.6\linewidth}}{\cellcolor[rgb]{0.9,0.9,0.9}{%
	\makecell[{{p{\linewidth}}}]{% 
	  \tt {\tiny [GM$\rangle$B]}  
	 Il tuo messaggio è stato inoltrato all'altro giocatore.\\ \tt \\ \tt \\ \tt You are currently playing round 3 of maximum 20.\\ \tt You have collectively accumulated 2 of 12 penalties.\\ \tt L'altro giocatore ha inviato questo messaggio:\\ \tt \\ \tt \\ \tt "Sì, confermo che il mio 'H' si trova nella prima riga e il mio 'T' si trova nella terza riga. Il mio 'W' si trova nella settima riga. Sono d'accordo con l'ordine standard di elencare gli oggetti in base all'alfabeto ('C', 'H', 'I', 'T', 'W'). Per procedere, potremmo iniziare a scambiare informazioni sulla posizione esatta di ogni oggetto, utilizzando le coordinate (x, y) per indicare la posizione di ogni oggetto. Ad esempio, potrei dire che il mio 'H' si trova nella posizione (3, 1). Come pensi di procedere?"\\ \tt \\ \tt \\ \tt Qual è il tuo prossimo comando? 
	  } 
	   } 
	   } 
	 & \\ 
 

    \theutterance \stepcounter{utterance}  

    & & & \multicolumn{4}{p{0.6\linewidth}}{\cellcolor[rgb]{1,0.85,0.72}{%
	\makecell[{{p{\linewidth}}}]{% 
	  \tt {\tiny [GM$\langle$B]}  
	 COMUNICA: Procediamo a scambiare informazioni sulla posizione esatta di ogni oggetto, utilizzando le coordinate (x, y) per indicare la posizione di ogni oggetto, partendo da 'C', poi 'H', 'I', 'T' e infine 'W', in modo da avere una visione chiara di dove sono posizionati tutti gli oggetti sui nostri tavoli. 
	  } 
	   } 
	   } 
	  \\ 
 

    \theutterance \stepcounter{utterance}  

    & & & \multicolumn{2}{p{0.3\linewidth}}{\cellcolor[rgb]{0.95,0.95,0.95}{%
	\makecell[{{p{\linewidth}}}]{% 
	  \tt {\tiny [GM$|$GM]}  
	 COMUNICA: Procediamo a scambiare informazioni sulla posizione esatta di ogni oggetto, utilizzando le coordinate (x, y) per indicare la posizione di ogni oggetto, partendo da 'C', poi 'H', 'I', 'T' e infine 'W', in modo da avere una visione chiara di dove sono posizionati tutti gli oggetti sui nostri tavoli. 
	  } 
	   } 
	   } 
	 & & \\ 
 

    \theutterance \stepcounter{utterance}  

    & & \multicolumn{4}{p{0.6\linewidth}}{\cellcolor[rgb]{0.9,0.9,0.9}{%
	\makecell[{{p{\linewidth}}}]{% 
	  \tt {\tiny [A$\langle$GM]}  
	 Il tuo messaggio è stato inoltrato all'altro giocatore.\\ \tt \\ \tt \\ \tt You are currently playing round 3 of maximum 20.\\ \tt You have collectively accumulated 2 of 12 penalties.\\ \tt L'altro giocatore ha inviato questo messaggio:\\ \tt \\ \tt \\ \tt "Procediamo a scambiare informazioni sulla posizione esatta di ogni oggetto, utilizzando le coordinate (x, y) per indicare la posizione di ogni oggetto, partendo da 'C', poi 'H', 'I', 'T' e infine 'W', in modo da avere una visione chiara di dove sono posizionati tutti gli oggetti sui nostri tavoli."\\ \tt \\ \tt \\ \tt Qual è il tuo prossimo comando? 
	  } 
	   } 
	   } 
	 & \\ 
 

    \theutterance \stepcounter{utterance}  

    & \multicolumn{4}{p{0.6\linewidth}}{\cellcolor[rgb]{0.8,1,0.9}{%
	\makecell[{{p{\linewidth}}}]{% 
	  \tt {\tiny [A$\rangle$GM]}  
	 COMUNICA: Inizierò con il mio 'C', che si trova nella sesta riga, non ho la cassa esatta, il mio 'H' si trova nella posizione (3, 1), il mio 'I' si trova nella quinta riga, il mio 'T' si trova nella posizione (2, 3) e il mio 'W' si trova nella settima riga. Per favore, fornisci le posizioni dei tuoi oggetti utilizzando lo stesso formato. 
	  } 
	   } 
	   } 
	 & & \\ 
 

    \theutterance \stepcounter{utterance}  

    & & & \multicolumn{2}{p{0.3\linewidth}}{\cellcolor[rgb]{0.95,0.95,0.95}{%
	\makecell[{{p{\linewidth}}}]{% 
	  \tt {\tiny [GM$|$GM]}  
	 COMUNICA: Inizierò con il mio 'C', che si trova nella sesta riga, non ho la cassa esatta, il mio 'H' si trova nella posizione (3, 1), il mio 'I' si trova nella quinta riga, il mio 'T' si trova nella posizione (2, 3) e il mio 'W' si trova nella settima riga. Per favore, fornisci le posizioni dei tuoi oggetti utilizzando lo stesso formato. 
	  } 
	   } 
	   } 
	 & & \\ 
 

    \theutterance \stepcounter{utterance}  

    & & \multicolumn{4}{p{0.6\linewidth}}{\cellcolor[rgb]{0.9,0.9,0.9}{%
	\makecell[{{p{\linewidth}}}]{% 
	  \tt {\tiny [GM$\rangle$B]}  
	 Il tuo messaggio è stato inoltrato all'altro giocatore.\\ \tt \\ \tt \\ \tt You are currently playing round 4 of maximum 20.\\ \tt You have collectively accumulated 2 of 12 penalties.\\ \tt L'altro giocatore ha inviato questo messaggio:\\ \tt \\ \tt \\ \tt "Inizierò con il mio 'C', che si trova nella sesta riga, non ho la cassa esatta, il mio 'H' si trova nella posizione (3, 1), il mio 'I' si trova nella quinta riga, il mio 'T' si trova nella posizione (2, 3) e il mio 'W' si trova nella settima riga. Per favore, fornisci le posizioni dei tuoi oggetti utilizzando lo stesso formato."\\ \tt \\ \tt \\ \tt Qual è il tuo prossimo comando? 
	  } 
	   } 
	   } 
	 & \\ 
 

    \theutterance \stepcounter{utterance}  

    & & & \multicolumn{4}{p{0.6\linewidth}}{\cellcolor[rgb]{1,0.85,0.72}{%
	\makecell[{{p{\linewidth}}}]{% 
	  \tt {\tiny [GM$\langle$B]}  
	 COMUNICA: Il mio 'C' si trova nella posizione (4, 4), il mio 'H' si trova nella posizione (1, 5), il mio 'I' si trova nella posizione (3, 4), il mio 'T' si trova nella posizione (3, 1) e il mio 'W' si trova nella posizione (1, 1). 
	  } 
	   } 
	   } 
	  \\ 
 

    \theutterance \stepcounter{utterance}  

    & & & \multicolumn{2}{p{0.3\linewidth}}{\cellcolor[rgb]{0.95,0.95,0.95}{%
	\makecell[{{p{\linewidth}}}]{% 
	  \tt {\tiny [GM$|$GM]}  
	 COMUNICA: Il mio 'C' si trova nella posizione (4, 4), il mio 'H' si trova nella posizione (1, 5), il mio 'I' si trova nella posizione (3, 4), il mio 'T' si trova nella posizione (3, 1) e il mio 'W' si trova nella posizione (1, 1). 
	  } 
	   } 
	   } 
	 & & \\ 
 

    \theutterance \stepcounter{utterance}  

    & & \multicolumn{4}{p{0.6\linewidth}}{\cellcolor[rgb]{0.9,0.9,0.9}{%
	\makecell[{{p{\linewidth}}}]{% 
	  \tt {\tiny [A$\langle$GM]}  
	 Il tuo messaggio è stato inoltrato all'altro giocatore.\\ \tt \\ \tt \\ \tt You are currently playing round 4 of maximum 20.\\ \tt You have collectively accumulated 2 of 12 penalties.\\ \tt L'altro giocatore ha inviato questo messaggio:\\ \tt \\ \tt \\ \tt "Il mio 'C' si trova nella posizione (4, 4), il mio 'H' si trova nella posizione (1, 5), il mio 'I' si trova nella posizione (3, 4), il mio 'T' si trova nella posizione (3, 1) e il mio 'W' si trova nella posizione (1, 1)."\\ \tt \\ \tt \\ \tt Qual è il tuo prossimo comando? 
	  } 
	   } 
	   } 
	 & \\ 
 

    \theutterance \stepcounter{utterance}  

    & \multicolumn{4}{p{0.6\linewidth}}{\cellcolor[rgb]{0.8,1,0.9}{%
	\makecell[{{p{\linewidth}}}]{% 
	  \tt {\tiny [A$\rangle$GM]}  
	 COMUNICA: Grazie per le informazioni! Ora che abbiamo le posizioni di tutti gli oggetti, possiamo iniziare a pianificare gli spostamenti. Notiamo che il mio 'H' e il tuo 'T' sono già nelle stesse righe, quindi potremmo iniziare a lavorare sugli oggetti che sono più lontani dalle loro posizioni di destinazione. Il mio 'W' è nella settima riga e il tuo 'W' è nella prima riga, quindi potremmo iniziare a spostare i 'W' per allinearli. Cosa ne pensi? Dovremmo spostare il mio 'W' verso l'alto o il tuo 'W' verso il basso? 
	  } 
	   } 
	   } 
	 & & \\ 
 

    \theutterance \stepcounter{utterance}  

    & & & \multicolumn{2}{p{0.3\linewidth}}{\cellcolor[rgb]{0.95,0.95,0.95}{%
	\makecell[{{p{\linewidth}}}]{% 
	  \tt {\tiny [GM$|$GM]}  
	 COMUNICA: Grazie per le informazioni! Ora che abbiamo le posizioni di tutti gli oggetti, possiamo iniziare a pianificare gli spostamenti. Notiamo che il mio 'H' e il tuo 'T' sono già nelle stesse righe, quindi potremmo iniziare a lavorare sugli oggetti che sono più lontani dalle loro posizioni di destinazione. Il mio 'W' è nella settima riga e il tuo 'W' è nella prima riga, quindi potremmo iniziare a spostare i 'W' per allinearli. Cosa ne pensi? Dovremmo spostare il mio 'W' verso l'alto o il tuo 'W' verso il basso? 
	  } 
	   } 
	   } 
	 & & \\ 
 

    \theutterance \stepcounter{utterance}  

    & & \multicolumn{4}{p{0.6\linewidth}}{\cellcolor[rgb]{0.9,0.9,0.9}{%
	\makecell[{{p{\linewidth}}}]{% 
	  \tt {\tiny [GM$\rangle$B]}  
	 Il tuo messaggio è stato inoltrato all'altro giocatore.\\ \tt \\ \tt \\ \tt You are currently playing round 5 of maximum 20.\\ \tt You have collectively accumulated 2 of 12 penalties.\\ \tt L'altro giocatore ha inviato questo messaggio:\\ \tt \\ \tt \\ \tt "Grazie per le informazioni! Ora che abbiamo le posizioni di tutti gli oggetti, possiamo iniziare a pianificare gli spostamenti. Notiamo che il mio 'H' e il tuo 'T' sono già nelle stesse righe, quindi potremmo iniziare a lavorare sugli oggetti che sono più lontani dalle loro posizioni di destinazione. Il mio 'W' è nella settima riga e il tuo 'W' è nella prima riga, quindi potremmo iniziare a spostare i 'W' per allinearli. Cosa ne pensi? Dovremmo spostare il mio 'W' verso l'alto o il tuo 'W' verso il basso?"\\ \tt \\ \tt \\ \tt Qual è il tuo prossimo comando? 
	  } 
	   } 
	   } 
	 & \\ 
 

    \theutterance \stepcounter{utterance}  

    & & & \multicolumn{4}{p{0.6\linewidth}}{\cellcolor[rgb]{1,0.85,0.72}{%
	\makecell[{{p{\linewidth}}}]{% 
	  \tt {\tiny [GM$\langle$B]}  
	 COMUNICA: Penso che sarebbe più efficiente spostare il tuo 'W' verso l'alto, poiché il mio 'W' è già nella prima riga e sarebbe più facile spostare il tuo 'W' verso l'alto rispetto a spostare il mio 'W' verso il basso, considerando le posizioni degli altri oggetti. 
	  } 
	   } 
	   } 
	  \\ 
 

    \theutterance \stepcounter{utterance}  

    & & & \multicolumn{2}{p{0.3\linewidth}}{\cellcolor[rgb]{0.95,0.95,0.95}{%
	\makecell[{{p{\linewidth}}}]{% 
	  \tt {\tiny [GM$|$GM]}  
	 COMUNICA: Penso che sarebbe più efficiente spostare il tuo 'W' verso l'alto, poiché il mio 'W' è già nella prima riga e sarebbe più facile spostare il tuo 'W' verso l'alto rispetto a spostare il mio 'W' verso il basso, considerando le posizioni degli altri oggetti. 
	  } 
	   } 
	   } 
	 & & \\ 
 

    \theutterance \stepcounter{utterance}  

    & & \multicolumn{4}{p{0.6\linewidth}}{\cellcolor[rgb]{0.9,0.9,0.9}{%
	\makecell[{{p{\linewidth}}}]{% 
	  \tt {\tiny [A$\langle$GM]}  
	 Il tuo messaggio è stato inoltrato all'altro giocatore.\\ \tt \\ \tt \\ \tt You are currently playing round 5 of maximum 20.\\ \tt You have collectively accumulated 2 of 12 penalties.\\ \tt L'altro giocatore ha inviato questo messaggio:\\ \tt \\ \tt \\ \tt "Penso che sarebbe più efficiente spostare il tuo 'W' verso l'alto, poiché il mio 'W' è già nella prima riga e sarebbe più facile spostare il tuo 'W' verso l'alto rispetto a spostare il mio 'W' verso il basso, considerando le posizioni degli altri oggetti."\\ \tt \\ \tt \\ \tt Qual è il tuo prossimo comando? 
	  } 
	   } 
	   } 
	 & \\ 
 

    \theutterance \stepcounter{utterance}  

    & \multicolumn{4}{p{0.6\linewidth}}{\cellcolor[rgb]{0.8,1,0.9}{%
	\makecell[{{p{\linewidth}}}]{% 
	  \tt {\tiny [A$\rangle$GM]}  
	 COMUNICA: Sono d'accordo con la tua proposta. Spostare il mio 'W' verso l'alto sembra la scelta più efficiente. Prima di fare la mossa, vorrei confermare che la posizione di destinazione per il mio 'W' sia la prima riga. Il mio 'W' si trova nella settima riga, quindi dovrò spostarlo di 6 posizioni verso l'alto. Cosa ne pensi? Dovrei procedere con lo spostamento del mio 'W'? 
	  } 
	   } 
	   } 
	 & & \\ 
 

    \theutterance \stepcounter{utterance}  

    & & & \multicolumn{2}{p{0.3\linewidth}}{\cellcolor[rgb]{0.95,0.95,0.95}{%
	\makecell[{{p{\linewidth}}}]{% 
	  \tt {\tiny [GM$|$GM]}  
	 COMUNICA: Sono d'accordo con la tua proposta. Spostare il mio 'W' verso l'alto sembra la scelta più efficiente. Prima di fare la mossa, vorrei confermare che la posizione di destinazione per il mio 'W' sia la prima riga. Il mio 'W' si trova nella settima riga, quindi dovrò spostarlo di 6 posizioni verso l'alto. Cosa ne pensi? Dovrei procedere con lo spostamento del mio 'W'? 
	  } 
	   } 
	   } 
	 & & \\ 
 

    \theutterance \stepcounter{utterance}  

    & & \multicolumn{4}{p{0.6\linewidth}}{\cellcolor[rgb]{0.9,0.9,0.9}{%
	\makecell[{{p{\linewidth}}}]{% 
	  \tt {\tiny [GM$\rangle$B]}  
	 Il tuo messaggio è stato inoltrato all'altro giocatore.\\ \tt \\ \tt \\ \tt You are currently playing round 6 of maximum 20.\\ \tt You have collectively accumulated 2 of 12 penalties.\\ \tt L'altro giocatore ha inviato questo messaggio:\\ \tt \\ \tt \\ \tt "Sono d'accordo con la tua proposta. Spostare il mio 'W' verso l'alto sembra la scelta più efficiente. Prima di fare la mossa, vorrei confermare che la posizione di destinazione per il mio 'W' sia la prima riga. Il mio 'W' si trova nella settima riga, quindi dovrò spostarlo di 6 posizioni verso l'alto. Cosa ne pensi? Dovrei procedere con lo spostamento del mio 'W'?"\\ \tt \\ \tt \\ \tt Qual è il tuo prossimo comando? 
	  } 
	   } 
	   } 
	 & \\ 
 

    \theutterance \stepcounter{utterance}  

    & & & \multicolumn{4}{p{0.6\linewidth}}{\cellcolor[rgb]{1,0.85,0.72}{%
	\makecell[{{p{\linewidth}}}]{% 
	  \tt {\tiny [GM$\langle$B]}  
	 COMUNICA: Sì, confermo che la posizione di destinazione per il tuo 'W' sia la prima riga, e quindi dovrebbe essere allineato con il mio 'W'. Spostare il tuo 'W' di 6 posizioni verso l'alto sembra la mossa giusta, quindi procedi pure con lo spostamento del tuo 'W' nella posizione (1,1) o nella colonna che abbiamo stabilito per il 'W'. 
	  } 
	   } 
	   } 
	  \\ 
 

    \theutterance \stepcounter{utterance}  

    & & & \multicolumn{2}{p{0.3\linewidth}}{\cellcolor[rgb]{0.95,0.95,0.95}{%
	\makecell[{{p{\linewidth}}}]{% 
	  \tt {\tiny [GM$|$GM]}  
	 COMUNICA: Sì, confermo che la posizione di destinazione per il tuo 'W' sia la prima riga, e quindi dovrebbe essere allineato con il mio 'W'. Spostare il tuo 'W' di 6 posizioni verso l'alto sembra la mossa giusta, quindi procedi pure con lo spostamento del tuo 'W' nella posizione (1,1) o nella colonna che abbiamo stabilito per il 'W'. 
	  } 
	   } 
	   } 
	 & & \\ 
 

    \theutterance \stepcounter{utterance}  

    & & \multicolumn{4}{p{0.6\linewidth}}{\cellcolor[rgb]{0.9,0.9,0.9}{%
	\makecell[{{p{\linewidth}}}]{% 
	  \tt {\tiny [A$\langle$GM]}  
	 Il tuo messaggio è stato inoltrato all'altro giocatore.\\ \tt \\ \tt \\ \tt You are currently playing round 6 of maximum 20.\\ \tt You have collectively accumulated 2 of 12 penalties.\\ \tt L'altro giocatore ha inviato questo messaggio:\\ \tt \\ \tt \\ \tt "Sì, confermo che la posizione di destinazione per il tuo 'W' sia la prima riga, e quindi dovrebbe essere allineato con il mio 'W'. Spostare il tuo 'W' di 6 posizioni verso l'alto sembra la mossa giusta, quindi procedi pure con lo spostamento del tuo 'W' nella posizione (1,1) o nella colonna che abbiamo stabilito per il 'W'."\\ \tt \\ \tt \\ \tt Qual è il tuo prossimo comando? 
	  } 
	   } 
	   } 
	 & \\ 
 

    \theutterance \stepcounter{utterance}  

    & \multicolumn{4}{p{0.6\linewidth}}{\cellcolor[rgb]{0.8,1,0.9}{%
	\makecell[{{p{\linewidth}}}]{% 
	  \tt {\tiny [A$\rangle$GM]}  
	 SPOSTA: W, (1, 1) 
	  } 
	   } 
	   } 
	 & & \\ 
 

    \theutterance \stepcounter{utterance}  

    & & & \multicolumn{2}{p{0.3\linewidth}}{\cellcolor[rgb]{0.95,0.95,0.95}{%
	\makecell[{{p{\linewidth}}}]{% 
	  \tt {\tiny [GM$|$GM]}  
	 SPOSTA: W, (1, 1) 
	  } 
	   } 
	   } 
	 & & \\ 
 

    \theutterance \stepcounter{utterance}  

    & & & \multicolumn{2}{p{0.3\linewidth}}{\cellcolor[rgb]{0.95,0.95,0.95}{%
	\makecell[{{p{\linewidth}}}]{% 
	  \tt {\tiny [GM$|$GM]}  
	 Spostato 'W' in (1,1) con successo. La tua griglia aggiornata appare così:\\ \tt ```\\ \tt  1234567\\ \tt ╔═╤═╤═╤═╗\\ \tt ║W│◌│H│◌║ 1\\ \tt ╟─┴─┼─┤◌║ 2\\ \tt ║◌T◌│◌└─╢ 3\\ \tt ╟───┘◌◌◌║ 4\\ \tt ║◌◌◌I◌┌─╢ 5\\ \tt ║◌◌◌C◌│◌║ 6\\ \tt ║◌◌◌◌◌│◌║ 7\\ \tt ╚═════╧═╝\\ \tt ```\\ \tt  
	  } 
	   } 
	   } 
	 & & \\ 
 

    \theutterance \stepcounter{utterance}  

    & & \multicolumn{4}{p{0.6\linewidth}}{\cellcolor[rgb]{0.9,0.9,0.9}{%
	\makecell[{{p{\linewidth}}}]{% 
	  \tt {\tiny [GM$\rangle$B]}  
	 Il tuo messaggio è stato inoltrato all'altro giocatore.\\ \tt \\ \tt \\ \tt You are currently playing round 7 of maximum 20.\\ \tt You have collectively accumulated 2 of 12 penalties.\\ \tt L'altro giocatore ha spostato un oggetto sulla sua griglia.\\ \tt \\ \tt \\ \tt Qual è il tuo prossimo comando? 
	  } 
	   } 
	   } 
	 & \\ 
 

    \theutterance \stepcounter{utterance}  

    & & & \multicolumn{4}{p{0.6\linewidth}}{\cellcolor[rgb]{1,0.85,0.72}{%
	\makecell[{{p{\linewidth}}}]{% 
	  \tt {\tiny [GM$\langle$B]}  
	 COMUNICA: Adesso che hai spostato il tuo 'W', potresti comunicarmi la sua nuova posizione esatta, utilizzando le coordinate (x, y), in modo che io possa aggiornare la mia visione della griglia e pianificare il prossimo spostamento? 
	  } 
	   } 
	   } 
	  \\ 
 

    \theutterance \stepcounter{utterance}  

    & & & \multicolumn{2}{p{0.3\linewidth}}{\cellcolor[rgb]{0.95,0.95,0.95}{%
	\makecell[{{p{\linewidth}}}]{% 
	  \tt {\tiny [GM$|$GM]}  
	 COMUNICA: Adesso che hai spostato il tuo 'W', potresti comunicarmi la sua nuova posizione esatta, utilizzando le coordinate (x, y), in modo che io possa aggiornare la mia visione della griglia e pianificare il prossimo spostamento? 
	  } 
	   } 
	   } 
	 & & \\ 
 

    \theutterance \stepcounter{utterance}  

    & & \multicolumn{4}{p{0.6\linewidth}}{\cellcolor[rgb]{0.9,0.9,0.9}{%
	\makecell[{{p{\linewidth}}}]{% 
	  \tt {\tiny [A$\langle$GM]}  
	 Spostato 'W' in (1,1) con successo. La tua griglia aggiornata appare così:\\ \tt ```\\ \tt  1234567\\ \tt ╔═╤═╤═╤═╗\\ \tt ║W│◌│H│◌║ 1\\ \tt ╟─┴─┼─┤◌║ 2\\ \tt ║◌T◌│◌└─╢ 3\\ \tt ╟───┘◌◌◌║ 4\\ \tt ║◌◌◌I◌┌─╢ 5\\ \tt ║◌◌◌C◌│◌║ 6\\ \tt ║◌◌◌◌◌│◌║ 7\\ \tt ╚═════╧═╝\\ \tt ```\\ \tt You are currently playing round 7 of maximum 20.\\ \tt You have collectively accumulated 2 of 12 penalties.\\ \tt L'altro giocatore ha inviato questo messaggio:\\ \tt \\ \tt \\ \tt "Adesso che hai spostato il tuo 'W', potresti comunicarmi la sua nuova posizione esatta, utilizzando le coordinate (x, y), in modo che io possa aggiornare la mia visione della griglia e pianificare il prossimo spostamento?"\\ \tt \\ \tt \\ \tt Qual è il tuo prossimo comando? 
	  } 
	   } 
	   } 
	 & \\ 
 

    \theutterance \stepcounter{utterance}  

    & \multicolumn{4}{p{0.6\linewidth}}{\cellcolor[rgb]{0.8,1,0.9}{%
	\makecell[{{p{\linewidth}}}]{% 
	  \tt {\tiny [A$\rangle$GM]}  
	 COMUNICA: Il mio 'W' è stato spostato con successo nella posizione (1, 1). La sua nuova posizione esatta è (1, 1). 
	  } 
	   } 
	   } 
	 & & \\ 
 

    \theutterance \stepcounter{utterance}  

    & & & \multicolumn{2}{p{0.3\linewidth}}{\cellcolor[rgb]{0.95,0.95,0.95}{%
	\makecell[{{p{\linewidth}}}]{% 
	  \tt {\tiny [GM$|$GM]}  
	 COMUNICA: Il mio 'W' è stato spostato con successo nella posizione (1, 1). La sua nuova posizione esatta è (1, 1). 
	  } 
	   } 
	   } 
	 & & \\ 
 

    \theutterance \stepcounter{utterance}  

    & & \multicolumn{4}{p{0.6\linewidth}}{\cellcolor[rgb]{0.9,0.9,0.9}{%
	\makecell[{{p{\linewidth}}}]{% 
	  \tt {\tiny [GM$\rangle$B]}  
	 Il tuo messaggio è stato inoltrato all'altro giocatore.\\ \tt \\ \tt \\ \tt You are currently playing round 8 of maximum 20.\\ \tt You have collectively accumulated 2 of 12 penalties.\\ \tt L'altro giocatore ha inviato questo messaggio:\\ \tt \\ \tt \\ \tt "Il mio 'W' è stato spostato con successo nella posizione (1, 1). La sua nuova posizione esatta è (1, 1)."\\ \tt \\ \tt \\ \tt Qual è il tuo prossimo comando? 
	  } 
	   } 
	   } 
	 & \\ 
 

    \theutterance \stepcounter{utterance}  

    & & & \multicolumn{4}{p{0.6\linewidth}}{\cellcolor[rgb]{1,0.85,0.72}{%
	\makecell[{{p{\linewidth}}}]{% 
	  \tt {\tiny [GM$\langle$B]}  
	 COMUNICA: Ora che il tuo 'W' è nella posizione (1, 1), potremmo passare al prossimo oggetto da allineare, che potrebbe essere il 'C' o l''I', a seconda di quale sia più vicino alla sua posizione di destinazione. Qual è la tua proposta per il prossimo passo? 
	  } 
	   } 
	   } 
	  \\ 
 

    \theutterance \stepcounter{utterance}  

    & & & \multicolumn{2}{p{0.3\linewidth}}{\cellcolor[rgb]{0.95,0.95,0.95}{%
	\makecell[{{p{\linewidth}}}]{% 
	  \tt {\tiny [GM$|$GM]}  
	 COMUNICA: Ora che il tuo 'W' è nella posizione (1, 1), potremmo passare al prossimo oggetto da allineare, che potrebbe essere il 'C' o l''I', a seconda di quale sia più vicino alla sua posizione di destinazione. Qual è la tua proposta per il prossimo passo? 
	  } 
	   } 
	   } 
	 & & \\ 
 

    \theutterance \stepcounter{utterance}  

    & & \multicolumn{4}{p{0.6\linewidth}}{\cellcolor[rgb]{0.9,0.9,0.9}{%
	\makecell[{{p{\linewidth}}}]{% 
	  \tt {\tiny [A$\langle$GM]}  
	 Il tuo messaggio è stato inoltrato all'altro giocatore.\\ \tt \\ \tt \\ \tt You are currently playing round 8 of maximum 20.\\ \tt You have collectively accumulated 2 of 12 penalties.\\ \tt L'altro giocatore ha inviato questo messaggio:\\ \tt \\ \tt \\ \tt "Ora che il tuo 'W' è nella posizione (1, 1), potremmo passare al prossimo oggetto da allineare, che potrebbe essere il 'C' o l''I', a seconda di quale sia più vicino alla sua posizione di destinazione. Qual è la tua proposta per il prossimo passo?"\\ \tt \\ \tt \\ \tt Qual è il tuo prossimo comando? 
	  } 
	   } 
	   } 
	 & \\ 
 

    \theutterance \stepcounter{utterance}  

    & \multicolumn{4}{p{0.6\linewidth}}{\cellcolor[rgb]{0.8,1,0.9}{%
	\makecell[{{p{\linewidth}}}]{% 
	  \tt {\tiny [A$\rangle$GM]}  
	 COMUNICA: Considerando le posizioni attuali, il mio 'C' è nella sesta riga e il tuo 'C' è nella posizione (4, 4). Il mio 'I' è nella quinta riga e il tuo 'I' è nella posizione (3, 4). Sembra che il mio 'I' sia più vicino alla sua posizione di destinazione, quindi potremmo iniziare a lavorare sull'allineamento dell''I'. Cosa ne pensi? Dovremmo spostare il mio 'I' per allineararlo con il tuo 'I'? 
	  } 
	   } 
	   } 
	 & & \\ 
 

    \theutterance \stepcounter{utterance}  

    & & & \multicolumn{2}{p{0.3\linewidth}}{\cellcolor[rgb]{0.95,0.95,0.95}{%
	\makecell[{{p{\linewidth}}}]{% 
	  \tt {\tiny [GM$|$GM]}  
	 COMUNICA: Considerando le posizioni attuali, il mio 'C' è nella sesta riga e il tuo 'C' è nella posizione (4, 4). Il mio 'I' è nella quinta riga e il tuo 'I' è nella posizione (3, 4). Sembra che il mio 'I' sia più vicino alla sua posizione di destinazione, quindi potremmo iniziare a lavorare sull'allineamento dell''I'. Cosa ne pensi? Dovremmo spostare il mio 'I' per allineararlo con il tuo 'I'? 
	  } 
	   } 
	   } 
	 & & \\ 
 

    \theutterance \stepcounter{utterance}  

    & & \multicolumn{4}{p{0.6\linewidth}}{\cellcolor[rgb]{0.9,0.9,0.9}{%
	\makecell[{{p{\linewidth}}}]{% 
	  \tt {\tiny [GM$\rangle$B]}  
	 Il tuo messaggio è stato inoltrato all'altro giocatore.\\ \tt \\ \tt \\ \tt You are currently playing round 9 of maximum 20.\\ \tt You have collectively accumulated 2 of 12 penalties.\\ \tt L'altro giocatore ha inviato questo messaggio:\\ \tt \\ \tt \\ \tt "Considerando le posizioni attuali, il mio 'C' è nella sesta riga e il tuo 'C' è nella posizione (4, 4). Il mio 'I' è nella quinta riga e il tuo 'I' è nella posizione (3, 4). Sembra che il mio 'I' sia più vicino alla sua posizione di destinazione, quindi potremmo iniziare a lavorare sull'allineamento dell''I'. Cosa ne pensi? Dovremmo spostare il mio 'I' per allineararlo con il tuo 'I'?"\\ \tt \\ \tt \\ \tt Qual è il tuo prossimo comando? 
	  } 
	   } 
	   } 
	 & \\ 
 

    \theutterance \stepcounter{utterance}  

    & & & \multicolumn{4}{p{0.6\linewidth}}{\cellcolor[rgb]{1,0.85,0.72}{%
	\makecell[{{p{\linewidth}}}]{% 
	  \tt {\tiny [GM$\langle$B]}  
	 COMUNICA: Sono d'accordo che il tuo 'I' sia più vicino alla sua posizione di destinazione. Tuttavia, considerando che il mio 'I' è nella posizione (3, 4) e il tuo 'I' è nella quinta riga, potremmo valutare se spostare il mio 'I' verso il basso o il tuo 'I' verso l'alto per allinearli. Qual è la tua valutazione sullo spostamento più efficiente? 
	  } 
	   } 
	   } 
	  \\ 
 

    \theutterance \stepcounter{utterance}  

    & & & \multicolumn{2}{p{0.3\linewidth}}{\cellcolor[rgb]{0.95,0.95,0.95}{%
	\makecell[{{p{\linewidth}}}]{% 
	  \tt {\tiny [GM$|$GM]}  
	 COMUNICA: Sono d'accordo che il tuo 'I' sia più vicino alla sua posizione di destinazione. Tuttavia, considerando che il mio 'I' è nella posizione (3, 4) e il tuo 'I' è nella quinta riga, potremmo valutare se spostare il mio 'I' verso il basso o il tuo 'I' verso l'alto per allinearli. Qual è la tua valutazione sullo spostamento più efficiente? 
	  } 
	   } 
	   } 
	 & & \\ 
 

    \theutterance \stepcounter{utterance}  

    & & \multicolumn{4}{p{0.6\linewidth}}{\cellcolor[rgb]{0.9,0.9,0.9}{%
	\makecell[{{p{\linewidth}}}]{% 
	  \tt {\tiny [A$\langle$GM]}  
	 Il tuo messaggio è stato inoltrato all'altro giocatore.\\ \tt \\ \tt \\ \tt You are currently playing round 9 of maximum 20.\\ \tt You have collectively accumulated 2 of 12 penalties.\\ \tt L'altro giocatore ha inviato questo messaggio:\\ \tt \\ \tt \\ \tt "Sono d'accordo che il tuo 'I' sia più vicino alla sua posizione di destinazione. Tuttavia, considerando che il mio 'I' è nella posizione (3, 4) e il tuo 'I' è nella quinta riga, potremmo valutare se spostare il mio 'I' verso il basso o il tuo 'I' verso l'alto per allinearli. Qual è la tua valutazione sullo spostamento più efficiente?"\\ \tt \\ \tt \\ \tt Qual è il tuo prossimo comando? 
	  } 
	   } 
	   } 
	 & \\ 
 

    \theutterance \stepcounter{utterance}  

    & \multicolumn{4}{p{0.6\linewidth}}{\cellcolor[rgb]{0.8,1,0.9}{%
	\makecell[{{p{\linewidth}}}]{% 
	  \tt {\tiny [A$\rangle$GM]}  
	 COMUNICA: Considerando le posizioni, sembra che spostare il mio 'I' verso l'alto nella posizione (3, 4) sarebbe più efficiente, poiché il mio 'I' è già nella colonna giusta, ma nella riga sbagliata. Spostare il mio 'I' verso l'alto di una posizione sembra la scelta più semplice. Cosa ne pensi? Dovrei procedere con lo spostamento del mio 'I'? 
	  } 
	   } 
	   } 
	 & & \\ 
 

    \theutterance \stepcounter{utterance}  

    & & & \multicolumn{2}{p{0.3\linewidth}}{\cellcolor[rgb]{0.95,0.95,0.95}{%
	\makecell[{{p{\linewidth}}}]{% 
	  \tt {\tiny [GM$|$GM]}  
	 COMUNICA: Considerando le posizioni, sembra che spostare il mio 'I' verso l'alto nella posizione (3, 4) sarebbe più efficiente, poiché il mio 'I' è già nella colonna giusta, ma nella riga sbagliata. Spostare il mio 'I' verso l'alto di una posizione sembra la scelta più semplice. Cosa ne pensi? Dovrei procedere con lo spostamento del mio 'I'? 
	  } 
	   } 
	   } 
	 & & \\ 
 

    \theutterance \stepcounter{utterance}  

    & & \multicolumn{4}{p{0.6\linewidth}}{\cellcolor[rgb]{0.9,0.9,0.9}{%
	\makecell[{{p{\linewidth}}}]{% 
	  \tt {\tiny [GM$\rangle$B]}  
	 Il tuo messaggio è stato inoltrato all'altro giocatore.\\ \tt \\ \tt \\ \tt You are currently playing round 10 of maximum 20.\\ \tt You have collectively accumulated 2 of 12 penalties.\\ \tt L'altro giocatore ha inviato questo messaggio:\\ \tt \\ \tt \\ \tt "Considerando le posizioni, sembra che spostare il mio 'I' verso l'alto nella posizione (3, 4) sarebbe più efficiente, poiché il mio 'I' è già nella colonna giusta, ma nella riga sbagliata. Spostare il mio 'I' verso l'alto di una posizione sembra la scelta più semplice. Cosa ne pensi? Dovrei procedere con lo spostamento del mio 'I'?"\\ \tt \\ \tt \\ \tt Qual è il tuo prossimo comando? 
	  } 
	   } 
	   } 
	 & \\ 
 

    \theutterance \stepcounter{utterance}  

    & & & \multicolumn{4}{p{0.6\linewidth}}{\cellcolor[rgb]{1,0.85,0.72}{%
	\makecell[{{p{\linewidth}}}]{% 
	  \tt {\tiny [GM$\langle$B]}  
	 COMUNICA: Sì, sono d'accordo che spostare il tuo 'I' verso l'alto nella posizione (3, 4) sembri la scelta più efficiente, considerando che è già nella colonna giusta. Procedi pure con lo spostamento del tuo 'I' nella posizione (3, 4) per allinearli. 
	  } 
	   } 
	   } 
	  \\ 
 

    \theutterance \stepcounter{utterance}  

    & & & \multicolumn{2}{p{0.3\linewidth}}{\cellcolor[rgb]{0.95,0.95,0.95}{%
	\makecell[{{p{\linewidth}}}]{% 
	  \tt {\tiny [GM$|$GM]}  
	 COMUNICA: Sì, sono d'accordo che spostare il tuo 'I' verso l'alto nella posizione (3, 4) sembri la scelta più efficiente, considerando che è già nella colonna giusta. Procedi pure con lo spostamento del tuo 'I' nella posizione (3, 4) per allinearli. 
	  } 
	   } 
	   } 
	 & & \\ 
 

    \theutterance \stepcounter{utterance}  

    & & \multicolumn{4}{p{0.6\linewidth}}{\cellcolor[rgb]{0.9,0.9,0.9}{%
	\makecell[{{p{\linewidth}}}]{% 
	  \tt {\tiny [A$\langle$GM]}  
	 Il tuo messaggio è stato inoltrato all'altro giocatore.\\ \tt \\ \tt \\ \tt You are currently playing round 10 of maximum 20.\\ \tt You have collectively accumulated 2 of 12 penalties.\\ \tt L'altro giocatore ha inviato questo messaggio:\\ \tt \\ \tt \\ \tt "Sì, sono d'accordo che spostare il tuo 'I' verso l'alto nella posizione (3, 4) sembri la scelta più efficiente, considerando che è già nella colonna giusta. Procedi pure con lo spostamento del tuo 'I' nella posizione (3, 4) per allinearli."\\ \tt \\ \tt \\ \tt Qual è il tuo prossimo comando? 
	  } 
	   } 
	   } 
	 & \\ 
 

    \theutterance \stepcounter{utterance}  

    & \multicolumn{4}{p{0.6\linewidth}}{\cellcolor[rgb]{0.8,1,0.9}{%
	\makecell[{{p{\linewidth}}}]{% 
	  \tt {\tiny [A$\rangle$GM]}  
	 SPOSTA: I, (3, 4) 
	  } 
	   } 
	   } 
	 & & \\ 
 

    \theutterance \stepcounter{utterance}  

    & & & \multicolumn{2}{p{0.3\linewidth}}{\cellcolor[rgb]{0.95,0.95,0.95}{%
	\makecell[{{p{\linewidth}}}]{% 
	  \tt {\tiny [GM$|$GM]}  
	 SPOSTA: I, (3, 4) 
	  } 
	   } 
	   } 
	 & & \\ 
 

    \theutterance \stepcounter{utterance}  

    & & & \multicolumn{2}{p{0.3\linewidth}}{\cellcolor[rgb]{0.95,0.95,0.95}{%
	\makecell[{{p{\linewidth}}}]{% 
	  \tt {\tiny [GM$|$GM]}  
	 Penalità: (3,4) non è vuoto, ma contiene '─'.\\ \tt \\ \tt You have collectively accumulated 3 of 12 penalties.\\ \tt Riprova! 
	  } 
	   } 
	   } 
	 & & \\ 
 

    \theutterance \stepcounter{utterance}  

    & & \multicolumn{4}{p{0.6\linewidth}}{\cellcolor[rgb]{0.9,0.9,0.9}{%
	\makecell[{{p{\linewidth}}}]{% 
	  \tt {\tiny [A$\langle$GM]}  
	 Penalità: (3,4) non è vuoto, ma contiene '─'.\\ \tt \\ \tt You have collectively accumulated 3 of 12 penalties.\\ \tt Riprova! 
	  } 
	   } 
	   } 
	 & \\ 
 

    \theutterance \stepcounter{utterance}  

    & \multicolumn{4}{p{0.6\linewidth}}{\cellcolor[rgb]{0.8,1,0.9}{%
	\makecell[{{p{\linewidth}}}]{% 
	  \tt {\tiny [A$\rangle$GM]}  
	 SPOSTA: I, (5, 4) 
	  } 
	   } 
	   } 
	 & & \\ 
 

    \theutterance \stepcounter{utterance}  

    & & & \multicolumn{2}{p{0.3\linewidth}}{\cellcolor[rgb]{0.95,0.95,0.95}{%
	\makecell[{{p{\linewidth}}}]{% 
	  \tt {\tiny [GM$|$GM]}  
	 SPOSTA: I, (5, 4) 
	  } 
	   } 
	   } 
	 & & \\ 
 

    \theutterance \stepcounter{utterance}  

    & & & \multicolumn{2}{p{0.3\linewidth}}{\cellcolor[rgb]{0.95,0.95,0.95}{%
	\makecell[{{p{\linewidth}}}]{% 
	  \tt {\tiny [GM$|$GM]}  
	 Spostato 'I' in (5,4) con successo. La tua griglia aggiornata appare così:\\ \tt ```\\ \tt  1234567\\ \tt ╔═╤═╤═╤═╗\\ \tt ║W│◌│H│◌║ 1\\ \tt ╟─┴─┼─┤◌║ 2\\ \tt ║◌T◌│◌└─╢ 3\\ \tt ╟───┘I◌◌║ 4\\ \tt ║◌◌◌◌◌┌─╢ 5\\ \tt ║◌◌◌C◌│◌║ 6\\ \tt ║◌◌◌◌◌│◌║ 7\\ \tt ╚═════╧═╝\\ \tt ```\\ \tt  
	  } 
	   } 
	   } 
	 & & \\ 
 

    \theutterance \stepcounter{utterance}  

    & & \multicolumn{4}{p{0.6\linewidth}}{\cellcolor[rgb]{0.9,0.9,0.9}{%
	\makecell[{{p{\linewidth}}}]{% 
	  \tt {\tiny [GM$\rangle$B]}  
	 Il tuo messaggio è stato inoltrato all'altro giocatore.\\ \tt \\ \tt \\ \tt You are currently playing round 11 of maximum 20.\\ \tt You have collectively accumulated 3 of 12 penalties.\\ \tt L'altro giocatore ha spostato un oggetto sulla sua griglia.\\ \tt \\ \tt \\ \tt Qual è il tuo prossimo comando? 
	  } 
	   } 
	   } 
	 & \\ 
 

    \theutterance \stepcounter{utterance}  

    & & & \multicolumn{4}{p{0.6\linewidth}}{\cellcolor[rgb]{1,0.85,0.72}{%
	\makecell[{{p{\linewidth}}}]{% 
	  \tt {\tiny [GM$\langle$B]}  
	 COMUNICA: Adesso che hai spostato il tuo 'I', potresti confermare se il tuo 'I' si trova ora nella posizione (3, 4) e se le posizioni degli altri oggetti sono cambiate? Questo mi aiuterà a valutare il prossimo passo. 
	  } 
	   } 
	   } 
	  \\ 
 

    \theutterance \stepcounter{utterance}  

    & & & \multicolumn{2}{p{0.3\linewidth}}{\cellcolor[rgb]{0.95,0.95,0.95}{%
	\makecell[{{p{\linewidth}}}]{% 
	  \tt {\tiny [GM$|$GM]}  
	 COMUNICA: Adesso che hai spostato il tuo 'I', potresti confermare se il tuo 'I' si trova ora nella posizione (3, 4) e se le posizioni degli altri oggetti sono cambiate? Questo mi aiuterà a valutare il prossimo passo. 
	  } 
	   } 
	   } 
	 & & \\ 
 

    \theutterance \stepcounter{utterance}  

    & & \multicolumn{4}{p{0.6\linewidth}}{\cellcolor[rgb]{0.9,0.9,0.9}{%
	\makecell[{{p{\linewidth}}}]{% 
	  \tt {\tiny [A$\langle$GM]}  
	 Spostato 'I' in (5,4) con successo. La tua griglia aggiornata appare così:\\ \tt ```\\ \tt  1234567\\ \tt ╔═╤═╤═╤═╗\\ \tt ║W│◌│H│◌║ 1\\ \tt ╟─┴─┼─┤◌║ 2\\ \tt ║◌T◌│◌└─╢ 3\\ \tt ╟───┘I◌◌║ 4\\ \tt ║◌◌◌◌◌┌─╢ 5\\ \tt ║◌◌◌C◌│◌║ 6\\ \tt ║◌◌◌◌◌│◌║ 7\\ \tt ╚═════╧═╝\\ \tt ```\\ \tt You are currently playing round 11 of maximum 20.\\ \tt You have collectively accumulated 3 of 12 penalties.\\ \tt L'altro giocatore ha inviato questo messaggio:\\ \tt \\ \tt \\ \tt "Adesso che hai spostato il tuo 'I', potresti confermare se il tuo 'I' si trova ora nella posizione (3, 4) e se le posizioni degli altri oggetti sono cambiate? Questo mi aiuterà a valutare il prossimo passo."\\ \tt \\ \tt \\ \tt Qual è il tuo prossimo comando? 
	  } 
	   } 
	   } 
	 & \\ 
 

    \theutterance \stepcounter{utterance}  

    & \multicolumn{4}{p{0.6\linewidth}}{\cellcolor[rgb]{0.8,1,0.9}{%
	\makecell[{{p{\linewidth}}}]{% 
	  \tt {\tiny [A$\rangle$GM]}  
	 COMUNICA: Il mio 'I' non si trova nella posizione (3, 4) come previsto, ma si trova nella posizione (5, 4) a causa dell'errore di destinazione. Le posizioni degli altri oggetti non sono cambiate, tranne per il fatto che il mio 'I' è stato spostato nella posizione (5, 4). Il mio 'W' è nella posizione (1, 1), il mio 'H' è nella posizione (3, 1), il mio 'T' è nella posizione (2, 3) e il mio 'C' è nella posizione (4, 6). 
	  } 
	   } 
	   } 
	 & & \\ 
 

    \theutterance \stepcounter{utterance}  

    & & & \multicolumn{2}{p{0.3\linewidth}}{\cellcolor[rgb]{0.95,0.95,0.95}{%
	\makecell[{{p{\linewidth}}}]{% 
	  \tt {\tiny [GM$|$GM]}  
	 COMUNICA: Il mio 'I' non si trova nella posizione (3, 4) come previsto, ma si trova nella posizione (5, 4) a causa dell'errore di destinazione. Le posizioni degli altri oggetti non sono cambiate, tranne per il fatto che il mio 'I' è stato spostato nella posizione (5, 4). Il mio 'W' è nella posizione (1, 1), il mio 'H' è nella posizione (3, 1), il mio 'T' è nella posizione (2, 3) e il mio 'C' è nella posizione (4, 6). 
	  } 
	   } 
	   } 
	 & & \\ 
 

    \theutterance \stepcounter{utterance}  

    & & \multicolumn{4}{p{0.6\linewidth}}{\cellcolor[rgb]{0.9,0.9,0.9}{%
	\makecell[{{p{\linewidth}}}]{% 
	  \tt {\tiny [GM$\rangle$B]}  
	 Il tuo messaggio è stato inoltrato all'altro giocatore.\\ \tt \\ \tt \\ \tt You are currently playing round 12 of maximum 20.\\ \tt You have collectively accumulated 3 of 12 penalties.\\ \tt L'altro giocatore ha inviato questo messaggio:\\ \tt \\ \tt \\ \tt "Il mio 'I' non si trova nella posizione (3, 4) come previsto, ma si trova nella posizione (5, 4) a causa dell'errore di destinazione. Le posizioni degli altri oggetti non sono cambiate, tranne per il fatto che il mio 'I' è stato spostato nella posizione (5, 4). Il mio 'W' è nella posizione (1, 1), il mio 'H' è nella posizione (3, 1), il mio 'T' è nella posizione (2, 3) e il mio 'C' è nella posizione (4, 6)."\\ \tt \\ \tt \\ \tt Qual è il tuo prossimo comando? 
	  } 
	   } 
	   } 
	 & \\ 
 

    \theutterance \stepcounter{utterance}  

    & & & \multicolumn{4}{p{0.6\linewidth}}{\cellcolor[rgb]{1,0.85,0.72}{%
	\makecell[{{p{\linewidth}}}]{% 
	  \tt {\tiny [GM$\langle$B]}  
	 COMUNICA: Grazie per la correzione. Considerando che il tuo 'I' è stato spostato nella posizione (5, 4) e non nella posizione prevista (3, 4), dobbiamo rivedere il nostro piano. Il mio 'I' è ancora nella posizione (3, 4). Potremmo valutare se spostare il mio 'I' verso il basso per allinearli con il tuo 'I' nella posizione (5, 4), o se ci sono altre mosse più prioritarie da fare. Cosa suggerisci di fare adesso? 
	  } 
	   } 
	   } 
	  \\ 
 

    \theutterance \stepcounter{utterance}  

    & & & \multicolumn{2}{p{0.3\linewidth}}{\cellcolor[rgb]{0.95,0.95,0.95}{%
	\makecell[{{p{\linewidth}}}]{% 
	  \tt {\tiny [GM$|$GM]}  
	 COMUNICA: Grazie per la correzione. Considerando che il tuo 'I' è stato spostato nella posizione (5, 4) e non nella posizione prevista (3, 4), dobbiamo rivedere il nostro piano. Il mio 'I' è ancora nella posizione (3, 4). Potremmo valutare se spostare il mio 'I' verso il basso per allinearli con il tuo 'I' nella posizione (5, 4), o se ci sono altre mosse più prioritarie da fare. Cosa suggerisci di fare adesso? 
	  } 
	   } 
	   } 
	 & & \\ 
 

    \theutterance \stepcounter{utterance}  

    & & \multicolumn{4}{p{0.6\linewidth}}{\cellcolor[rgb]{0.9,0.9,0.9}{%
	\makecell[{{p{\linewidth}}}]{% 
	  \tt {\tiny [A$\langle$GM]}  
	 Il tuo messaggio è stato inoltrato all'altro giocatore.\\ \tt \\ \tt \\ \tt You are currently playing round 12 of maximum 20.\\ \tt You have collectively accumulated 3 of 12 penalties.\\ \tt L'altro giocatore ha inviato questo messaggio:\\ \tt \\ \tt \\ \tt "Grazie per la correzione. Considerando che il tuo 'I' è stato spostato nella posizione (5, 4) e non nella posizione prevista (3, 4), dobbiamo rivedere il nostro piano. Il mio 'I' è ancora nella posizione (3, 4). Potremmo valutare se spostare il mio 'I' verso il basso per allinearli con il tuo 'I' nella posizione (5, 4), o se ci sono altre mosse più prioritarie da fare. Cosa suggerisci di fare adesso?"\\ \tt \\ \tt \\ \tt Qual è il tuo prossimo comando? 
	  } 
	   } 
	   } 
	 & \\ 
 

    \theutterance \stepcounter{utterance}  

    & \multicolumn{4}{p{0.6\linewidth}}{\cellcolor[rgb]{0.8,1,0.9}{%
	\makecell[{{p{\linewidth}}}]{% 
	  \tt {\tiny [A$\rangle$GM]}  
	 COMUNICA: Considerando le posizioni attuali, potremmo valutare se spostare il mio 'T' per allineararlo con il tuo 'T', o se ci sono altre mosse più prioritarie da fare. Il mio 'T' è nella posizione (2, 3) e il tuo 'T' è nella posizione (3, 1). Spostare il mio 'T' nella posizione (3, 1) potrebbe essere una buona mossa. Cosa ne pensi? Dovremmo procedere con lo spostamento del mio 'T'? 
	  } 
	   } 
	   } 
	 & & \\ 
 

    \theutterance \stepcounter{utterance}  

    & & & \multicolumn{2}{p{0.3\linewidth}}{\cellcolor[rgb]{0.95,0.95,0.95}{%
	\makecell[{{p{\linewidth}}}]{% 
	  \tt {\tiny [GM$|$GM]}  
	 COMUNICA: Considerando le posizioni attuali, potremmo valutare se spostare il mio 'T' per allineararlo con il tuo 'T', o se ci sono altre mosse più prioritarie da fare. Il mio 'T' è nella posizione (2, 3) e il tuo 'T' è nella posizione (3, 1). Spostare il mio 'T' nella posizione (3, 1) potrebbe essere una buona mossa. Cosa ne pensi? Dovremmo procedere con lo spostamento del mio 'T'? 
	  } 
	   } 
	   } 
	 & & \\ 
 

    \theutterance \stepcounter{utterance}  

    & & \multicolumn{4}{p{0.6\linewidth}}{\cellcolor[rgb]{0.9,0.9,0.9}{%
	\makecell[{{p{\linewidth}}}]{% 
	  \tt {\tiny [GM$\rangle$B]}  
	 Il tuo messaggio è stato inoltrato all'altro giocatore.\\ \tt \\ \tt \\ \tt You are currently playing round 13 of maximum 20.\\ \tt You have collectively accumulated 3 of 12 penalties.\\ \tt L'altro giocatore ha inviato questo messaggio:\\ \tt \\ \tt \\ \tt "Considerando le posizioni attuali, potremmo valutare se spostare il mio 'T' per allineararlo con il tuo 'T', o se ci sono altre mosse più prioritarie da fare. Il mio 'T' è nella posizione (2, 3) e il tuo 'T' è nella posizione (3, 1). Spostare il mio 'T' nella posizione (3, 1) potrebbe essere una buona mossa. Cosa ne pensi? Dovremmo procedere con lo spostamento del mio 'T'?"\\ \tt \\ \tt \\ \tt Qual è il tuo prossimo comando? 
	  } 
	   } 
	   } 
	 & \\ 
 

    \theutterance \stepcounter{utterance}  

    & & & \multicolumn{4}{p{0.6\linewidth}}{\cellcolor[rgb]{1,0.85,0.72}{%
	\makecell[{{p{\linewidth}}}]{% 
	  \tt {\tiny [GM$\langle$B]}  
	 COMUNICA: Spostare il tuo 'T' nella posizione (3, 1) potrebbe essere una buona mossa, poiché allineerebbe i 'T' e semplificherebbe la griglia. Tuttavia, considerando che il mio 'T' è già nella posizione (3, 1), potremmo dover spostare il mio 'T' per fare spazio al tuo 'T'. Potremmo valutare se spostare il mio 'T' in una posizione vicina per poi spostare il tuo 'T' nella posizione (3, 1), o se c'è un'altra strategia più efficiente. Cosa ne pensi? 
	  } 
	   } 
	   } 
	  \\ 
 

    \theutterance \stepcounter{utterance}  

    & & & \multicolumn{2}{p{0.3\linewidth}}{\cellcolor[rgb]{0.95,0.95,0.95}{%
	\makecell[{{p{\linewidth}}}]{% 
	  \tt {\tiny [GM$|$GM]}  
	 COMUNICA: Spostare il tuo 'T' nella posizione (3, 1) potrebbe essere una buona mossa, poiché allineerebbe i 'T' e semplificherebbe la griglia. Tuttavia, considerando che il mio 'T' è già nella posizione (3, 1), potremmo dover spostare il mio 'T' per fare spazio al tuo 'T'. Potremmo valutare se spostare il mio 'T' in una posizione vicina per poi spostare il tuo 'T' nella posizione (3, 1), o se c'è un'altra strategia più efficiente. Cosa ne pensi? 
	  } 
	   } 
	   } 
	 & & \\ 
 

    \theutterance \stepcounter{utterance}  

    & & \multicolumn{4}{p{0.6\linewidth}}{\cellcolor[rgb]{0.9,0.9,0.9}{%
	\makecell[{{p{\linewidth}}}]{% 
	  \tt {\tiny [A$\langle$GM]}  
	 Il tuo messaggio è stato inoltrato all'altro giocatore.\\ \tt \\ \tt \\ \tt You are currently playing round 13 of maximum 20.\\ \tt You have collectively accumulated 3 of 12 penalties.\\ \tt L'altro giocatore ha inviato questo messaggio:\\ \tt \\ \tt \\ \tt "Spostare il tuo 'T' nella posizione (3, 1) potrebbe essere una buona mossa, poiché allineerebbe i 'T' e semplificherebbe la griglia. Tuttavia, considerando che il mio 'T' è già nella posizione (3, 1), potremmo dover spostare il mio 'T' per fare spazio al tuo 'T'. Potremmo valutare se spostare il mio 'T' in una posizione vicina per poi spostare il tuo 'T' nella posizione (3, 1), o se c'è un'altra strategia più efficiente. Cosa ne pensi?"\\ \tt \\ \tt \\ \tt Qual è il tuo prossimo comando? 
	  } 
	   } 
	   } 
	 & \\ 
 

    \theutterance \stepcounter{utterance}  

    & \multicolumn{4}{p{0.6\linewidth}}{\cellcolor[rgb]{0.8,1,0.9}{%
	\makecell[{{p{\linewidth}}}]{% 
	  \tt {\tiny [A$\rangle$GM]}  
	 COMUNICA: Sono d'accordo che spostare il mio 'T' nella posizione (3, 1) potrebbe essere complicato a causa della presenza del tuo 'T'. Una possibile strategia potrebbe essere di spostare il mio 'T' nella posizione (3, 2) o (3, 3) se possibile, e poi valutare il prossimo passo. Tuttavia, considerando la griglia attuale, potremmo anche valutare di spostare il mio 'H' per fare spazio o allineare gli altri oggetti. Cosa ne pensi di spostare il mio 'H' per semplificare la griglia? 
	  } 
	   } 
	   } 
	 & & \\ 
 

    \theutterance \stepcounter{utterance}  

    & & & \multicolumn{2}{p{0.3\linewidth}}{\cellcolor[rgb]{0.95,0.95,0.95}{%
	\makecell[{{p{\linewidth}}}]{% 
	  \tt {\tiny [GM$|$GM]}  
	 COMUNICA: Sono d'accordo che spostare il mio 'T' nella posizione (3, 1) potrebbe essere complicato a causa della presenza del tuo 'T'. Una possibile strategia potrebbe essere di spostare il mio 'T' nella posizione (3, 2) o (3, 3) se possibile, e poi valutare il prossimo passo. Tuttavia, considerando la griglia attuale, potremmo anche valutare di spostare il mio 'H' per fare spazio o allineare gli altri oggetti. Cosa ne pensi di spostare il mio 'H' per semplificare la griglia? 
	  } 
	   } 
	   } 
	 & & \\ 
 

    \theutterance \stepcounter{utterance}  

    & & \multicolumn{4}{p{0.6\linewidth}}{\cellcolor[rgb]{0.9,0.9,0.9}{%
	\makecell[{{p{\linewidth}}}]{% 
	  \tt {\tiny [GM$\rangle$B]}  
	 Il tuo messaggio è stato inoltrato all'altro giocatore.\\ \tt \\ \tt \\ \tt You are currently playing round 14 of maximum 20.\\ \tt You have collectively accumulated 3 of 12 penalties.\\ \tt L'altro giocatore ha inviato questo messaggio:\\ \tt \\ \tt \\ \tt "Sono d'accordo che spostare il mio 'T' nella posizione (3, 1) potrebbe essere complicato a causa della presenza del tuo 'T'. Una possibile strategia potrebbe essere di spostare il mio 'T' nella posizione (3, 2) o (3, 3) se possibile, e poi valutare il prossimo passo. Tuttavia, considerando la griglia attuale, potremmo anche valutare di spostare il mio 'H' per fare spazio o allineare gli altri oggetti. Cosa ne pensi di spostare il mio 'H' per semplificare la griglia?"\\ \tt \\ \tt \\ \tt Qual è il tuo prossimo comando? 
	  } 
	   } 
	   } 
	 & \\ 
 

    \theutterance \stepcounter{utterance}  

    & & & \multicolumn{4}{p{0.6\linewidth}}{\cellcolor[rgb]{1,0.85,0.72}{%
	\makecell[{{p{\linewidth}}}]{% 
	  \tt {\tiny [GM$\langle$B]}  
	 COMUNICA: Spostare il tuo 'H' potrebbe essere un'ottima mossa per semplificare la griglia e creare spazio per altri oggetti. Considerando che il mio 'H' è nella posizione (1, 5) e il tuo 'H' è nella posizione (3, 1), potremmo valutare se spostare il tuo 'H' nella posizione (1, 1) o in una posizione vicina al mio 'H' per allinearli. Questo potrebbe aiutare a ridurre la complessità della griglia e facilitare gli spostamenti futuri. Cosa ne pensi di spostare il tuo 'H' in una di queste posizioni? 
	  } 
	   } 
	   } 
	  \\ 
 

    \theutterance \stepcounter{utterance}  

    & & & \multicolumn{2}{p{0.3\linewidth}}{\cellcolor[rgb]{0.95,0.95,0.95}{%
	\makecell[{{p{\linewidth}}}]{% 
	  \tt {\tiny [GM$|$GM]}  
	 COMUNICA: Spostare il tuo 'H' potrebbe essere un'ottima mossa per semplificare la griglia e creare spazio per altri oggetti. Considerando che il mio 'H' è nella posizione (1, 5) e il tuo 'H' è nella posizione (3, 1), potremmo valutare se spostare il tuo 'H' nella posizione (1, 1) o in una posizione vicina al mio 'H' per allinearli. Questo potrebbe aiutare a ridurre la complessità della griglia e facilitare gli spostamenti futuri. Cosa ne pensi di spostare il tuo 'H' in una di queste posizioni? 
	  } 
	   } 
	   } 
	 & & \\ 
 

    \theutterance \stepcounter{utterance}  

    & & \multicolumn{4}{p{0.6\linewidth}}{\cellcolor[rgb]{0.9,0.9,0.9}{%
	\makecell[{{p{\linewidth}}}]{% 
	  \tt {\tiny [A$\langle$GM]}  
	 Il tuo messaggio è stato inoltrato all'altro giocatore.\\ \tt \\ \tt \\ \tt You are currently playing round 14 of maximum 20.\\ \tt You have collectively accumulated 3 of 12 penalties.\\ \tt L'altro giocatore ha inviato questo messaggio:\\ \tt \\ \tt \\ \tt "Spostare il tuo 'H' potrebbe essere un'ottima mossa per semplificare la griglia e creare spazio per altri oggetti. Considerando che il mio 'H' è nella posizione (1, 5) e il tuo 'H' è nella posizione (3, 1), potremmo valutare se spostare il tuo 'H' nella posizione (1, 1) o in una posizione vicina al mio 'H' per allinearli. Questo potrebbe aiutare a ridurre la complessità della griglia e facilitare gli spostamenti futuri. Cosa ne pensi di spostare il tuo 'H' in una di queste posizioni?"\\ \tt \\ \tt \\ \tt Qual è il tuo prossimo comando? 
	  } 
	   } 
	   } 
	 & \\ 
 

    \theutterance \stepcounter{utterance}  

    & \multicolumn{4}{p{0.6\linewidth}}{\cellcolor[rgb]{0.8,1,0.9}{%
	\makecell[{{p{\linewidth}}}]{% 
	  \tt {\tiny [A$\rangle$GM]}  
	 COMUNICA: Non posso spostare il mio 'H' nella posizione (1, 1) perché è già occupata dal mio 'W'. Tuttavia, potremmo valutare di spostare il mio 'H' in una posizione vicina al mio 'W', come ad esempio (2, 1), in modo da liberare spazio nella colonna 3 e semplificare la griglia. Cosa ne pensi? Dovremmo procedere con lo spostamento del mio 'H' in una posizione vicina al mio 'W'? 
	  } 
	   } 
	   } 
	 & & \\ 
 

    \theutterance \stepcounter{utterance}  

    & & & \multicolumn{2}{p{0.3\linewidth}}{\cellcolor[rgb]{0.95,0.95,0.95}{%
	\makecell[{{p{\linewidth}}}]{% 
	  \tt {\tiny [GM$|$GM]}  
	 COMUNICA: Non posso spostare il mio 'H' nella posizione (1, 1) perché è già occupata dal mio 'W'. Tuttavia, potremmo valutare di spostare il mio 'H' in una posizione vicina al mio 'W', come ad esempio (2, 1), in modo da liberare spazio nella colonna 3 e semplificare la griglia. Cosa ne pensi? Dovremmo procedere con lo spostamento del mio 'H' in una posizione vicina al mio 'W'? 
	  } 
	   } 
	   } 
	 & & \\ 
 

    \theutterance \stepcounter{utterance}  

    & & \multicolumn{4}{p{0.6\linewidth}}{\cellcolor[rgb]{0.9,0.9,0.9}{%
	\makecell[{{p{\linewidth}}}]{% 
	  \tt {\tiny [GM$\rangle$B]}  
	 Il tuo messaggio è stato inoltrato all'altro giocatore.\\ \tt \\ \tt \\ \tt You are currently playing round 15 of maximum 20.\\ \tt You have collectively accumulated 3 of 12 penalties.\\ \tt L'altro giocatore ha inviato questo messaggio:\\ \tt \\ \tt \\ \tt "Non posso spostare il mio 'H' nella posizione (1, 1) perché è già occupata dal mio 'W'. Tuttavia, potremmo valutare di spostare il mio 'H' in una posizione vicina al mio 'W', come ad esempio (2, 1), in modo da liberare spazio nella colonna 3 e semplificare la griglia. Cosa ne pensi? Dovremmo procedere con lo spostamento del mio 'H' in una posizione vicina al mio 'W'?"\\ \tt \\ \tt \\ \tt Qual è il tuo prossimo comando? 
	  } 
	   } 
	   } 
	 & \\ 
 

    \theutterance \stepcounter{utterance}  

    & & & \multicolumn{4}{p{0.6\linewidth}}{\cellcolor[rgb]{1,0.85,0.72}{%
	\makecell[{{p{\linewidth}}}]{% 
	  \tt {\tiny [GM$\langle$B]}  
	 COMUNICA: Sì, sono d'accordo che spostare il tuo 'H' nella posizione (2, 1) potrebbe essere una buona mossa per liberare spazio nella colonna 3 e semplificare la griglia. Questo potrebbe anche aiutare a creare più opzioni per gli spostamenti futuri degli altri oggetti. Procedi pure con lo spostamento del tuo 'H' nella posizione (2, 1). 
	  } 
	   } 
	   } 
	  \\ 
 

    \theutterance \stepcounter{utterance}  

    & & & \multicolumn{2}{p{0.3\linewidth}}{\cellcolor[rgb]{0.95,0.95,0.95}{%
	\makecell[{{p{\linewidth}}}]{% 
	  \tt {\tiny [GM$|$GM]}  
	 COMUNICA: Sì, sono d'accordo che spostare il tuo 'H' nella posizione (2, 1) potrebbe essere una buona mossa per liberare spazio nella colonna 3 e semplificare la griglia. Questo potrebbe anche aiutare a creare più opzioni per gli spostamenti futuri degli altri oggetti. Procedi pure con lo spostamento del tuo 'H' nella posizione (2, 1). 
	  } 
	   } 
	   } 
	 & & \\ 
 

    \theutterance \stepcounter{utterance}  

    & & \multicolumn{4}{p{0.6\linewidth}}{\cellcolor[rgb]{0.9,0.9,0.9}{%
	\makecell[{{p{\linewidth}}}]{% 
	  \tt {\tiny [A$\langle$GM]}  
	 Il tuo messaggio è stato inoltrato all'altro giocatore.\\ \tt \\ \tt \\ \tt You are currently playing round 15 of maximum 20.\\ \tt You have collectively accumulated 3 of 12 penalties.\\ \tt L'altro giocatore ha inviato questo messaggio:\\ \tt \\ \tt \\ \tt "Sì, sono d'accordo che spostare il tuo 'H' nella posizione (2, 1) potrebbe essere una buona mossa per liberare spazio nella colonna 3 e semplificare la griglia. Questo potrebbe anche aiutare a creare più opzioni per gli spostamenti futuri degli altri oggetti. Procedi pure con lo spostamento del tuo 'H' nella posizione (2, 1)."\\ \tt \\ \tt \\ \tt Qual è il tuo prossimo comando? 
	  } 
	   } 
	   } 
	 & \\ 
 

    \theutterance \stepcounter{utterance}  

    & \multicolumn{4}{p{0.6\linewidth}}{\cellcolor[rgb]{0.8,1,0.9}{%
	\makecell[{{p{\linewidth}}}]{% 
	  \tt {\tiny [A$\rangle$GM]}  
	 SPOSTA: H, (2, 1) 
	  } 
	   } 
	   } 
	 & & \\ 
 

    \theutterance \stepcounter{utterance}  

    & & & \multicolumn{2}{p{0.3\linewidth}}{\cellcolor[rgb]{0.95,0.95,0.95}{%
	\makecell[{{p{\linewidth}}}]{% 
	  \tt {\tiny [GM$|$GM]}  
	 SPOSTA: H, (2, 1) 
	  } 
	   } 
	   } 
	 & & \\ 
 

    \theutterance \stepcounter{utterance}  

    & & & \multicolumn{2}{p{0.3\linewidth}}{\cellcolor[rgb]{0.95,0.95,0.95}{%
	\makecell[{{p{\linewidth}}}]{% 
	  \tt {\tiny [GM$|$GM]}  
	 Penalità: (2,1) non è vuoto, ma contiene '│'.\\ \tt \\ \tt You have collectively accumulated 4 of 12 penalties.\\ \tt Riprova! 
	  } 
	   } 
	   } 
	 & & \\ 
 

    \theutterance \stepcounter{utterance}  

    & & \multicolumn{4}{p{0.6\linewidth}}{\cellcolor[rgb]{0.9,0.9,0.9}{%
	\makecell[{{p{\linewidth}}}]{% 
	  \tt {\tiny [A$\langle$GM]}  
	 Penalità: (2,1) non è vuoto, ma contiene '│'.\\ \tt \\ \tt You have collectively accumulated 4 of 12 penalties.\\ \tt Riprova! 
	  } 
	   } 
	   } 
	 & \\ 
 

    \theutterance \stepcounter{utterance}  

    & \multicolumn{4}{p{0.6\linewidth}}{\cellcolor[rgb]{0.8,1,0.9}{%
	\makecell[{{p{\linewidth}}}]{% 
	  \tt {\tiny [A$\rangle$GM]}  
	 SPOSTA: H, (2, 2) 
	  } 
	   } 
	   } 
	 & & \\ 
 

    \theutterance \stepcounter{utterance}  

    & & & \multicolumn{2}{p{0.3\linewidth}}{\cellcolor[rgb]{0.95,0.95,0.95}{%
	\makecell[{{p{\linewidth}}}]{% 
	  \tt {\tiny [GM$|$GM]}  
	 SPOSTA: H, (2, 2) 
	  } 
	   } 
	   } 
	 & & \\ 
 

    \theutterance \stepcounter{utterance}  

    & & & \multicolumn{2}{p{0.3\linewidth}}{\cellcolor[rgb]{0.95,0.95,0.95}{%
	\makecell[{{p{\linewidth}}}]{% 
	  \tt {\tiny [GM$|$GM]}  
	 Penalità: (2,2) non è vuoto, ma contiene '┴'.\\ \tt \\ \tt You have collectively accumulated 5 of 12 penalties.\\ \tt Riprova! 
	  } 
	   } 
	   } 
	 & & \\ 
 

    \theutterance \stepcounter{utterance}  

    & & \multicolumn{4}{p{0.6\linewidth}}{\cellcolor[rgb]{0.9,0.9,0.9}{%
	\makecell[{{p{\linewidth}}}]{% 
	  \tt {\tiny [A$\langle$GM]}  
	 Penalità: (2,2) non è vuoto, ma contiene '┴'.\\ \tt \\ \tt You have collectively accumulated 5 of 12 penalties.\\ \tt Riprova! 
	  } 
	   } 
	   } 
	 & \\ 
 

    \theutterance \stepcounter{utterance}  

    & \multicolumn{4}{p{0.6\linewidth}}{\cellcolor[rgb]{0.8,1,0.9}{%
	\makecell[{{p{\linewidth}}}]{% 
	  \tt {\tiny [A$\rangle$GM]}  
	 SPOSTA: H, (4, 1) 
	  } 
	   } 
	   } 
	 & & \\ 
 

    \theutterance \stepcounter{utterance}  

    & & & \multicolumn{2}{p{0.3\linewidth}}{\cellcolor[rgb]{0.95,0.95,0.95}{%
	\makecell[{{p{\linewidth}}}]{% 
	  \tt {\tiny [GM$|$GM]}  
	 SPOSTA: H, (4, 1) 
	  } 
	   } 
	   } 
	 & & \\ 
 

    \theutterance \stepcounter{utterance}  

    & & & \multicolumn{2}{p{0.3\linewidth}}{\cellcolor[rgb]{0.95,0.95,0.95}{%
	\makecell[{{p{\linewidth}}}]{% 
	  \tt {\tiny [GM$|$GM]}  
	 Penalità: (4,1) non è vuoto, ma contiene '│'.\\ \tt \\ \tt You have collectively accumulated 6 of 12 penalties.\\ \tt Riprova! 
	  } 
	   } 
	   } 
	 & & \\ 
 

    \theutterance \stepcounter{utterance}  

    & & \multicolumn{4}{p{0.6\linewidth}}{\cellcolor[rgb]{0.9,0.9,0.9}{%
	\makecell[{{p{\linewidth}}}]{% 
	  \tt {\tiny [A$\langle$GM]}  
	 Penalità: (4,1) non è vuoto, ma contiene '│'.\\ \tt \\ \tt You have collectively accumulated 6 of 12 penalties.\\ \tt Riprova! 
	  } 
	   } 
	   } 
	 & \\ 
 

    \theutterance \stepcounter{utterance}  

    & \multicolumn{4}{p{0.6\linewidth}}{\cellcolor[rgb]{0.8,1,0.9}{%
	\makecell[{{p{\linewidth}}}]{% 
	  \tt {\tiny [A$\rangle$GM]}  
	 SPOSTA: H, (5, 1) 
	  } 
	   } 
	   } 
	 & & \\ 
 

    \theutterance \stepcounter{utterance}  

    & & & \multicolumn{2}{p{0.3\linewidth}}{\cellcolor[rgb]{0.95,0.95,0.95}{%
	\makecell[{{p{\linewidth}}}]{% 
	  \tt {\tiny [GM$|$GM]}  
	 SPOSTA: H, (5, 1) 
	  } 
	   } 
	   } 
	 & & \\ 
 

    \theutterance \stepcounter{utterance}  

    & & & \multicolumn{2}{p{0.3\linewidth}}{\cellcolor[rgb]{0.95,0.95,0.95}{%
	\makecell[{{p{\linewidth}}}]{% 
	  \tt {\tiny [GM$|$GM]}  
	 Penalità: (5,1) non è vuoto, ma contiene 'H'.\\ \tt \\ \tt You have collectively accumulated 7 of 12 penalties.\\ \tt Riprova! 
	  } 
	   } 
	   } 
	 & & \\ 
 

    \theutterance \stepcounter{utterance}  

    & & \multicolumn{4}{p{0.6\linewidth}}{\cellcolor[rgb]{0.9,0.9,0.9}{%
	\makecell[{{p{\linewidth}}}]{% 
	  \tt {\tiny [A$\langle$GM]}  
	 Penalità: (5,1) non è vuoto, ma contiene 'H'.\\ \tt \\ \tt You have collectively accumulated 7 of 12 penalties.\\ \tt Riprova! 
	  } 
	   } 
	   } 
	 & \\ 
 

    \theutterance \stepcounter{utterance}  

    & \multicolumn{4}{p{0.6\linewidth}}{\cellcolor[rgb]{0.8,1,0.9}{%
	\makecell[{{p{\linewidth}}}]{% 
	  \tt {\tiny [A$\rangle$GM]}  
	 SPOSTA: H, (4, 2) 
	  } 
	   } 
	   } 
	 & & \\ 
 

    \theutterance \stepcounter{utterance}  

    & & & \multicolumn{2}{p{0.3\linewidth}}{\cellcolor[rgb]{0.95,0.95,0.95}{%
	\makecell[{{p{\linewidth}}}]{% 
	  \tt {\tiny [GM$|$GM]}  
	 SPOSTA: H, (4, 2) 
	  } 
	   } 
	   } 
	 & & \\ 
 

    \theutterance \stepcounter{utterance}  

    & & & \multicolumn{2}{p{0.3\linewidth}}{\cellcolor[rgb]{0.95,0.95,0.95}{%
	\makecell[{{p{\linewidth}}}]{% 
	  \tt {\tiny [GM$|$GM]}  
	 Penalità: (4,2) non è vuoto, ma contiene '┼'.\\ \tt \\ \tt You have collectively accumulated 8 of 12 penalties.\\ \tt Riprova! 
	  } 
	   } 
	   } 
	 & & \\ 
 

    \theutterance \stepcounter{utterance}  

    & & \multicolumn{4}{p{0.6\linewidth}}{\cellcolor[rgb]{0.9,0.9,0.9}{%
	\makecell[{{p{\linewidth}}}]{% 
	  \tt {\tiny [A$\langle$GM]}  
	 Penalità: (4,2) non è vuoto, ma contiene '┼'.\\ \tt \\ \tt You have collectively accumulated 8 of 12 penalties.\\ \tt Riprova! 
	  } 
	   } 
	   } 
	 & \\ 
 

    \theutterance \stepcounter{utterance}  

    & \multicolumn{4}{p{0.6\linewidth}}{\cellcolor[rgb]{0.8,1,0.9}{%
	\makecell[{{p{\linewidth}}}]{% 
	  \tt {\tiny [A$\rangle$GM]}  
	 SPOSTA: H, (5, 2) 
	  } 
	   } 
	   } 
	 & & \\ 
 

    \theutterance \stepcounter{utterance}  

    & & & \multicolumn{2}{p{0.3\linewidth}}{\cellcolor[rgb]{0.95,0.95,0.95}{%
	\makecell[{{p{\linewidth}}}]{% 
	  \tt {\tiny [GM$|$GM]}  
	 SPOSTA: H, (5, 2) 
	  } 
	   } 
	   } 
	 & & \\ 
 

    \theutterance \stepcounter{utterance}  

    & & & \multicolumn{2}{p{0.3\linewidth}}{\cellcolor[rgb]{0.95,0.95,0.95}{%
	\makecell[{{p{\linewidth}}}]{% 
	  \tt {\tiny [GM$|$GM]}  
	 Penalità: (5,2) non è vuoto, ma contiene '─'.\\ \tt \\ \tt You have collectively accumulated 9 of 12 penalties.\\ \tt Riprova! 
	  } 
	   } 
	   } 
	 & & \\ 
 

    \theutterance \stepcounter{utterance}  

    & & \multicolumn{4}{p{0.6\linewidth}}{\cellcolor[rgb]{0.9,0.9,0.9}{%
	\makecell[{{p{\linewidth}}}]{% 
	  \tt {\tiny [A$\langle$GM]}  
	 Penalità: (5,2) non è vuoto, ma contiene '─'.\\ \tt \\ \tt You have collectively accumulated 9 of 12 penalties.\\ \tt Riprova! 
	  } 
	   } 
	   } 
	 & \\ 
 

    \theutterance \stepcounter{utterance}  

    & \multicolumn{4}{p{0.6\linewidth}}{\cellcolor[rgb]{0.8,1,0.9}{%
	\makecell[{{p{\linewidth}}}]{% 
	  \tt {\tiny [A$\rangle$GM]}  
	 SPOSTA: H, (4, 3) 
	  } 
	   } 
	   } 
	 & & \\ 
 

    \theutterance \stepcounter{utterance}  

    & & & \multicolumn{2}{p{0.3\linewidth}}{\cellcolor[rgb]{0.95,0.95,0.95}{%
	\makecell[{{p{\linewidth}}}]{% 
	  \tt {\tiny [GM$|$GM]}  
	 SPOSTA: H, (4, 3) 
	  } 
	   } 
	   } 
	 & & \\ 
 

    \theutterance \stepcounter{utterance}  

    & & & \multicolumn{2}{p{0.3\linewidth}}{\cellcolor[rgb]{0.95,0.95,0.95}{%
	\makecell[{{p{\linewidth}}}]{% 
	  \tt {\tiny [GM$|$GM]}  
	 Penalità: (4,3) non è vuoto, ma contiene '│'.\\ \tt \\ \tt You have collectively accumulated 10 of 12 penalties.\\ \tt Riprova! 
	  } 
	   } 
	   } 
	 & & \\ 
 

    \theutterance \stepcounter{utterance}  

    & & \multicolumn{4}{p{0.6\linewidth}}{\cellcolor[rgb]{0.9,0.9,0.9}{%
	\makecell[{{p{\linewidth}}}]{% 
	  \tt {\tiny [A$\langle$GM]}  
	 Penalità: (4,3) non è vuoto, ma contiene '│'.\\ \tt \\ \tt You have collectively accumulated 10 of 12 penalties.\\ \tt Riprova! 
	  } 
	   } 
	   } 
	 & \\ 
 

    \theutterance \stepcounter{utterance}  

    & \multicolumn{4}{p{0.6\linewidth}}{\cellcolor[rgb]{0.8,1,0.9}{%
	\makecell[{{p{\linewidth}}}]{% 
	  \tt {\tiny [A$\rangle$GM]}  
	 SPOSTA: H, (4, 5) 
	  } 
	   } 
	   } 
	 & & \\ 
 

    \theutterance \stepcounter{utterance}  

    & & & \multicolumn{2}{p{0.3\linewidth}}{\cellcolor[rgb]{0.95,0.95,0.95}{%
	\makecell[{{p{\linewidth}}}]{% 
	  \tt {\tiny [GM$|$GM]}  
	 SPOSTA: H, (4, 5) 
	  } 
	   } 
	   } 
	 & & \\ 
 

    \theutterance \stepcounter{utterance}  

    & & & \multicolumn{2}{p{0.3\linewidth}}{\cellcolor[rgb]{0.95,0.95,0.95}{%
	\makecell[{{p{\linewidth}}}]{% 
	  \tt {\tiny [GM$|$GM]}  
	 Spostato 'H' in (4,5) con successo. La tua griglia aggiornata appare così:\\ \tt ```\\ \tt  1234567\\ \tt ╔═╤═╤═╤═╗\\ \tt ║W│◌│◌│◌║ 1\\ \tt ╟─┴─┼─┤◌║ 2\\ \tt ║◌T◌│◌└─╢ 3\\ \tt ╟───┘I◌◌║ 4\\ \tt ║◌◌◌H◌┌─╢ 5\\ \tt ║◌◌◌C◌│◌║ 6\\ \tt ║◌◌◌◌◌│◌║ 7\\ \tt ╚═════╧═╝\\ \tt ```\\ \tt  
	  } 
	   } 
	   } 
	 & & \\ 
 

    \theutterance \stepcounter{utterance}  

    & & \multicolumn{4}{p{0.6\linewidth}}{\cellcolor[rgb]{0.9,0.9,0.9}{%
	\makecell[{{p{\linewidth}}}]{% 
	  \tt {\tiny [GM$\rangle$B]}  
	 Il tuo messaggio è stato inoltrato all'altro giocatore.\\ \tt \\ \tt \\ \tt You are currently playing round 16 of maximum 20.\\ \tt You have collectively accumulated 10 of 12 penalties.\\ \tt L'altro giocatore ha spostato un oggetto sulla sua griglia.\\ \tt \\ \tt \\ \tt Qual è il tuo prossimo comando? 
	  } 
	   } 
	   } 
	 & \\ 
 

    \theutterance \stepcounter{utterance}  

    & & & \multicolumn{4}{p{0.6\linewidth}}{\cellcolor[rgb]{1,0.85,0.72}{%
	\makecell[{{p{\linewidth}}}]{% 
	  \tt {\tiny [GM$\langle$B]}  
	 COMUNICA: Adesso che hai spostato il tuo 'H', potresti comunicarmi la nuova posizione esatta di tutti gli oggetti sulla tua griglia, in modo che io possa aggiornare la mia visione della griglia e pianificare il prossimo spostamento? Questo mi aiuterà a capire come procedere per allineare gli oggetti. 
	  } 
	   } 
	   } 
	  \\ 
 

    \theutterance \stepcounter{utterance}  

    & & & \multicolumn{2}{p{0.3\linewidth}}{\cellcolor[rgb]{0.95,0.95,0.95}{%
	\makecell[{{p{\linewidth}}}]{% 
	  \tt {\tiny [GM$|$GM]}  
	 COMUNICA: Adesso che hai spostato il tuo 'H', potresti comunicarmi la nuova posizione esatta di tutti gli oggetti sulla tua griglia, in modo che io possa aggiornare la mia visione della griglia e pianificare il prossimo spostamento? Questo mi aiuterà a capire come procedere per allineare gli oggetti. 
	  } 
	   } 
	   } 
	 & & \\ 
 

    \theutterance \stepcounter{utterance}  

    & & \multicolumn{4}{p{0.6\linewidth}}{\cellcolor[rgb]{0.9,0.9,0.9}{%
	\makecell[{{p{\linewidth}}}]{% 
	  \tt {\tiny [A$\langle$GM]}  
	 Spostato 'H' in (4,5) con successo. La tua griglia aggiornata appare così:\\ \tt ```\\ \tt  1234567\\ \tt ╔═╤═╤═╤═╗\\ \tt ║W│◌│◌│◌║ 1\\ \tt ╟─┴─┼─┤◌║ 2\\ \tt ║◌T◌│◌└─╢ 3\\ \tt ╟───┘I◌◌║ 4\\ \tt ║◌◌◌H◌┌─╢ 5\\ \tt ║◌◌◌C◌│◌║ 6\\ \tt ║◌◌◌◌◌│◌║ 7\\ \tt ╚═════╧═╝\\ \tt ```\\ \tt You are currently playing round 16 of maximum 20.\\ \tt You have collectively accumulated 10 of 12 penalties.\\ \tt L'altro giocatore ha inviato questo messaggio:\\ \tt \\ \tt \\ \tt "Adesso che hai spostato il tuo 'H', potresti comunicarmi la nuova posizione esatta di tutti gli oggetti sulla tua griglia, in modo che io possa aggiornare la mia visione della griglia e pianificare il prossimo spostamento? Questo mi aiuterà a capire come procedere per allineare gli oggetti."\\ \tt \\ \tt \\ \tt Qual è il tuo prossimo comando? 
	  } 
	   } 
	   } 
	 & \\ 
 

    \theutterance \stepcounter{utterance}  

    & \multicolumn{4}{p{0.6\linewidth}}{\cellcolor[rgb]{0.8,1,0.9}{%
	\makecell[{{p{\linewidth}}}]{% 
	  \tt {\tiny [A$\rangle$GM]}  
	 COMUNICA: Ecco la nuova posizione esatta di tutti gli oggetti sulla mia griglia: il mio 'W' è nella posizione (1, 1), il mio 'T' è nella posizione (2, 3), il mio 'I' è nella posizione (5, 4), il mio 'C' è nella posizione (4, 6) e il mio 'H' è nella posizione (4, 5). Spero che questo ti aiuti a pianificare il prossimo spostamento. Cosa pensi di fare adesso? 
	  } 
	   } 
	   } 
	 & & \\ 
 

    \theutterance \stepcounter{utterance}  

    & & & \multicolumn{2}{p{0.3\linewidth}}{\cellcolor[rgb]{0.95,0.95,0.95}{%
	\makecell[{{p{\linewidth}}}]{% 
	  \tt {\tiny [GM$|$GM]}  
	 COMUNICA: Ecco la nuova posizione esatta di tutti gli oggetti sulla mia griglia: il mio 'W' è nella posizione (1, 1), il mio 'T' è nella posizione (2, 3), il mio 'I' è nella posizione (5, 4), il mio 'C' è nella posizione (4, 6) e il mio 'H' è nella posizione (4, 5). Spero che questo ti aiuti a pianificare il prossimo spostamento. Cosa pensi di fare adesso? 
	  } 
	   } 
	   } 
	 & & \\ 
 

    \theutterance \stepcounter{utterance}  

    & & \multicolumn{4}{p{0.6\linewidth}}{\cellcolor[rgb]{0.9,0.9,0.9}{%
	\makecell[{{p{\linewidth}}}]{% 
	  \tt {\tiny [GM$\rangle$B]}  
	 Il tuo messaggio è stato inoltrato all'altro giocatore.\\ \tt \\ \tt \\ \tt You are currently playing round 17 of maximum 20.\\ \tt You have collectively accumulated 10 of 12 penalties.\\ \tt L'altro giocatore ha inviato questo messaggio:\\ \tt \\ \tt \\ \tt "Ecco la nuova posizione esatta di tutti gli oggetti sulla mia griglia: il mio 'W' è nella posizione (1, 1), il mio 'T' è nella posizione (2, 3), il mio 'I' è nella posizione (5, 4), il mio 'C' è nella posizione (4, 6) e il mio 'H' è nella posizione (4, 5). Spero che questo ti aiuti a pianificare il prossimo spostamento. Cosa pensi di fare adesso?"\\ \tt \\ \tt \\ \tt Qual è il tuo prossimo comando? 
	  } 
	   } 
	   } 
	 & \\ 
 

    \theutterance \stepcounter{utterance}  

    & & & \multicolumn{4}{p{0.6\linewidth}}{\cellcolor[rgb]{1,0.85,0.72}{%
	\makecell[{{p{\linewidth}}}]{% 
	  \tt {\tiny [GM$\langle$B]}  
	 COMUNICA: Grazie per l'aggiornamento. Considerando le nuove posizioni, potremmo valutare di allineare gli oggetti che sono più vicini alle loro posizioni di destinazione. Notiamo che il mio 'C' è nella posizione (4, 4) e il tuo 'C' è nella posizione (4, 6), quindi potremmo cercare di allinearli. Tuttavia, considerando il numero di penalità accumulate, potremmo anche valutare se è il momento di dichiarare la fine del gioco. Cosa ne pensi? Dovremmo continuare a spostare gli oggetti o dichiarare la fine del gioco? 
	  } 
	   } 
	   } 
	  \\ 
 

    \theutterance \stepcounter{utterance}  

    & & & \multicolumn{2}{p{0.3\linewidth}}{\cellcolor[rgb]{0.95,0.95,0.95}{%
	\makecell[{{p{\linewidth}}}]{% 
	  \tt {\tiny [GM$|$GM]}  
	 COMUNICA: Grazie per l'aggiornamento. Considerando le nuove posizioni, potremmo valutare di allineare gli oggetti che sono più vicini alle loro posizioni di destinazione. Notiamo che il mio 'C' è nella posizione (4, 4) e il tuo 'C' è nella posizione (4, 6), quindi potremmo cercare di allinearli. Tuttavia, considerando il numero di penalità accumulate, potremmo anche valutare se è il momento di dichiarare la fine del gioco. Cosa ne pensi? Dovremmo continuare a spostare gli oggetti o dichiarare la fine del gioco? 
	  } 
	   } 
	   } 
	 & & \\ 
 

    \theutterance \stepcounter{utterance}  

    & & \multicolumn{4}{p{0.6\linewidth}}{\cellcolor[rgb]{0.9,0.9,0.9}{%
	\makecell[{{p{\linewidth}}}]{% 
	  \tt {\tiny [A$\langle$GM]}  
	 Il tuo messaggio è stato inoltrato all'altro giocatore.\\ \tt \\ \tt \\ \tt You are currently playing round 17 of maximum 20.\\ \tt You have collectively accumulated 10 of 12 penalties.\\ \tt L'altro giocatore ha inviato questo messaggio:\\ \tt \\ \tt \\ \tt "Grazie per l'aggiornamento. Considerando le nuove posizioni, potremmo valutare di allineare gli oggetti che sono più vicini alle loro posizioni di destinazione. Notiamo che il mio 'C' è nella posizione (4, 4) e il tuo 'C' è nella posizione (4, 6), quindi potremmo cercare di allinearli. Tuttavia, considerando il numero di penalità accumulate, potremmo anche valutare se è il momento di dichiarare la fine del gioco. Cosa ne pensi? Dovremmo continuare a spostare gli oggetti o dichiarare la fine del gioco?"\\ \tt \\ \tt \\ \tt Qual è il tuo prossimo comando? 
	  } 
	   } 
	   } 
	 & \\ 
 

    \theutterance \stepcounter{utterance}  

    & \multicolumn{4}{p{0.6\linewidth}}{\cellcolor[rgb]{0.8,1,0.9}{%
	\makecell[{{p{\linewidth}}}]{% 
	  \tt {\tiny [A$\rangle$GM]}  
	 COMUNICA: Sono d'accordo che il numero di penalità accumulate è alto e che potremmo valutare di dichiarare la fine del gioco. Tuttavia, considerando che siamo vicini all'obiettivo di allineare gli oggetti, potremmo cercare di fare ancora alcuni spostamenti strategici per migliorare la nostra posizione. Se non riusciamo a fare progressi significativi nei prossimi turni, potremmo considerare di dichiarare la fine del gioco. Cosa ne pensi di proporre un limite di turni per continuare a giocare? Se non miglioriamo la nostra posizione entro un certo numero di turni, potremmo dichiarare la fine del gioco. 
	  } 
	   } 
	   } 
	 & & \\ 
 

    \theutterance \stepcounter{utterance}  

    & & & \multicolumn{2}{p{0.3\linewidth}}{\cellcolor[rgb]{0.95,0.95,0.95}{%
	\makecell[{{p{\linewidth}}}]{% 
	  \tt {\tiny [GM$|$GM]}  
	 COMUNICA: Sono d'accordo che il numero di penalità accumulate è alto e che potremmo valutare di dichiarare la fine del gioco. Tuttavia, considerando che siamo vicini all'obiettivo di allineare gli oggetti, potremmo cercare di fare ancora alcuni spostamenti strategici per migliorare la nostra posizione. Se non riusciamo a fare progressi significativi nei prossimi turni, potremmo considerare di dichiarare la fine del gioco. Cosa ne pensi di proporre un limite di turni per continuare a giocare? Se non miglioriamo la nostra posizione entro un certo numero di turni, potremmo dichiarare la fine del gioco. 
	  } 
	   } 
	   } 
	 & & \\ 
 

    \theutterance \stepcounter{utterance}  

    & & \multicolumn{4}{p{0.6\linewidth}}{\cellcolor[rgb]{0.9,0.9,0.9}{%
	\makecell[{{p{\linewidth}}}]{% 
	  \tt {\tiny [GM$\rangle$B]}  
	 Il tuo messaggio è stato inoltrato all'altro giocatore.\\ \tt \\ \tt \\ \tt You are currently playing round 18 of maximum 20.\\ \tt You have collectively accumulated 10 of 12 penalties.\\ \tt L'altro giocatore ha inviato questo messaggio:\\ \tt \\ \tt \\ \tt "Sono d'accordo che il numero di penalità accumulate è alto e che potremmo valutare di dichiarare la fine del gioco. Tuttavia, considerando che siamo vicini all'obiettivo di allineare gli oggetti, potremmo cercare di fare ancora alcuni spostamenti strategici per migliorare la nostra posizione. Se non riusciamo a fare progressi significativi nei prossimi turni, potremmo considerare di dichiarare la fine del gioco. Cosa ne pensi di proporre un limite di turni per continuare a giocare? Se non miglioriamo la nostra posizione entro un certo numero di turni, potremmo dichiarare la fine del gioco."\\ \tt \\ \tt \\ \tt Qual è il tuo prossimo comando? 
	  } 
	   } 
	   } 
	 & \\ 
 

    \theutterance \stepcounter{utterance}  

    & & & \multicolumn{4}{p{0.6\linewidth}}{\cellcolor[rgb]{1,0.85,0.72}{%
	\makecell[{{p{\linewidth}}}]{% 
	  \tt {\tiny [GM$\langle$B]}  
	 COMUNICA: Sono d'accordo con la tua proposta di stabilire un limite di turni per continuare a giocare. Considerando che siamo già al round 18 e abbiamo solo 2 turni rimanenti prima di raggiungere il limite massimo di 20 round, potremmo stabilire che se non riusciamo a migliorare significativamente la nostra posizione entro i prossimi 2 turni, dichiariamo la fine del gioco. Questo ci darà una linea guida chiara per decidere quando fermarci. Sei d'accordo con questo piano? 
	  } 
	   } 
	   } 
	  \\ 
 

    \theutterance \stepcounter{utterance}  

    & & & \multicolumn{2}{p{0.3\linewidth}}{\cellcolor[rgb]{0.95,0.95,0.95}{%
	\makecell[{{p{\linewidth}}}]{% 
	  \tt {\tiny [GM$|$GM]}  
	 COMUNICA: Sono d'accordo con la tua proposta di stabilire un limite di turni per continuare a giocare. Considerando che siamo già al round 18 e abbiamo solo 2 turni rimanenti prima di raggiungere il limite massimo di 20 round, potremmo stabilire che se non riusciamo a migliorare significativamente la nostra posizione entro i prossimi 2 turni, dichiariamo la fine del gioco. Questo ci darà una linea guida chiara per decidere quando fermarci. Sei d'accordo con questo piano? 
	  } 
	   } 
	   } 
	 & & \\ 
 

    \theutterance \stepcounter{utterance}  

    & & \multicolumn{4}{p{0.6\linewidth}}{\cellcolor[rgb]{0.9,0.9,0.9}{%
	\makecell[{{p{\linewidth}}}]{% 
	  \tt {\tiny [A$\langle$GM]}  
	 Il tuo messaggio è stato inoltrato all'altro giocatore.\\ \tt \\ \tt \\ \tt You are currently playing round 18 of maximum 20.\\ \tt You have collectively accumulated 10 of 12 penalties.\\ \tt L'altro giocatore ha inviato questo messaggio:\\ \tt \\ \tt \\ \tt "Sono d'accordo con la tua proposta di stabilire un limite di turni per continuare a giocare. Considerando che siamo già al round 18 e abbiamo solo 2 turni rimanenti prima di raggiungere il limite massimo di 20 round, potremmo stabilire che se non riusciamo a migliorare significativamente la nostra posizione entro i prossimi 2 turni, dichiariamo la fine del gioco. Questo ci darà una linea guida chiara per decidere quando fermarci. Sei d'accordo con questo piano?"\\ \tt \\ \tt \\ \tt Qual è il tuo prossimo comando? 
	  } 
	   } 
	   } 
	 & \\ 
 

    \theutterance \stepcounter{utterance}  

    & \multicolumn{4}{p{0.6\linewidth}}{\cellcolor[rgb]{0.8,1,0.9}{%
	\makecell[{{p{\linewidth}}}]{% 
	  \tt {\tiny [A$\rangle$GM]}  
	 COMUNICA: Sono d'accordo con il tuo piano di stabilire un limite di 2 turni per continuare a giocare. Se non riusciamo a migliorare significativamente la nostra posizione entro i prossimi 2 turni, dichiariamo la fine del gioco. Sembra una strategia ragionevole considerando il numero di penalità accumulate e il tempo rimanente. Propongo di procedere con il prossimo turno e vedere se possiamo fare progressi. Se non miglioriamo, potremo valutare di dichiarare la fine del gioco al turno successivo. COMUNICA: finito? 
	  } 
	   } 
	   } 
	 & & \\ 
 

    \theutterance \stepcounter{utterance}  

    & & & \multicolumn{2}{p{0.3\linewidth}}{\cellcolor[rgb]{0.95,0.95,0.95}{%
	\makecell[{{p{\linewidth}}}]{% 
	  \tt {\tiny [GM$|$GM]}  
	 COMUNICA: Sono d'accordo con il tuo piano di stabilire un limite di 2 turni per continuare a giocare. Se non riusciamo a migliorare significativamente la nostra posizione entro i prossimi 2 turni, dichiariamo la fine del gioco. Sembra una strategia ragionevole considerando il numero di penalità accumulate e il tempo rimanente. Propongo di procedere con il prossimo turno e vedere se possiamo fare progressi. Se non miglioriamo, potremo valutare di dichiarare la fine del gioco al turno successivo. COMUNICA: finito? 
	  } 
	   } 
	   } 
	 & & \\ 
 

    \theutterance \stepcounter{utterance}  

    & & & \multicolumn{2}{p{0.3\linewidth}}{\cellcolor[rgb]{0.95,0.95,0.95}{%
	\makecell[{{p{\linewidth}}}]{% 
	  \tt {\tiny [GM$|$GM]}  
	 Invalid format: head is not empty: 'COMUNICA: Sono d'accordo con il tuo piano di stabilire un limite di 2 turni per continuare a giocare. Se non riusciamo a migliorare significativamente la nostra posizione entro i prossimi 2 turni, dichiariamo la fine del gioco. Sembra una strategia ragionevole considerando il numero di penalità accumulate e il tempo rimanente. Propongo di procedere con il prossimo turno e vedere se possiamo fare progressi. Se non miglioriamo, potremo valutare di dichiarare la fine del gioco al turno successivo. ' 
	  } 
	   } 
	   } 
	 & & \\ 
 

    \theutterance \stepcounter{utterance}  

    & & \multicolumn{4}{p{0.6\linewidth}}{\cellcolor[rgb]{0.9,0.9,0.9}{%
	\makecell[{{p{\linewidth}}}]{% 
	  \tt {\tiny [A$\langle$GM]}  
	 Penalità: Il tuo messaggio non deve contenere nulla prima del comando!\\ \tt \\ \tt Assicurati che la tua risposta contenga solo `COMUNICA: <MESSAGE>` o `SPOSTA: <OBJECT>, (<X>, <Y>)`, e nient'altro!\\ \tt \\ \tt You have collectively accumulated 11 of 12 penalties.\\ \tt Riprova! 
	  } 
	   } 
	   } 
	 & \\ 
 

    \theutterance \stepcounter{utterance}  

    & \multicolumn{4}{p{0.6\linewidth}}{\cellcolor[rgb]{0.8,1,0.9}{%
	\makecell[{{p{\linewidth}}}]{% 
	  \tt {\tiny [A$\rangle$GM]}  
	 COMUNICA: Sono d'accordo con il tuo piano di stabilire un limite di 2 turni per continuare a giocare. Se non riusciamo a migliorare significativamente la nostra posizione entro i prossimi 2 turni, dichiariamo la fine del gioco. Sembra una strategia ragionevole considerando il numero di penalità accumulate e il tempo rimanente. Propongo di procedere con il prossimo turno e vedere se possiamo fare progressi. Se non miglioriamo, potremo valutare di dichiarare la fine del gioco al turno successivo. COMUNICA: finito? 
	  } 
	   } 
	   } 
	 & & \\ 
 

    \theutterance \stepcounter{utterance}  

    & & & \multicolumn{2}{p{0.3\linewidth}}{\cellcolor[rgb]{0.95,0.95,0.95}{%
	\makecell[{{p{\linewidth}}}]{% 
	  \tt {\tiny [GM$|$GM]}  
	 COMUNICA: Sono d'accordo con il tuo piano di stabilire un limite di 2 turni per continuare a giocare. Se non riusciamo a migliorare significativamente la nostra posizione entro i prossimi 2 turni, dichiariamo la fine del gioco. Sembra una strategia ragionevole considerando il numero di penalità accumulate e il tempo rimanente. Propongo di procedere con il prossimo turno e vedere se possiamo fare progressi. Se non miglioriamo, potremo valutare di dichiarare la fine del gioco al turno successivo. COMUNICA: finito? 
	  } 
	   } 
	   } 
	 & & \\ 
 

    \theutterance \stepcounter{utterance}  

    & & & \multicolumn{2}{p{0.3\linewidth}}{\cellcolor[rgb]{0.95,0.95,0.95}{%
	\makecell[{{p{\linewidth}}}]{% 
	  \tt {\tiny [GM$|$GM]}  
	 Invalid format: head is not empty: 'COMUNICA: Sono d'accordo con il tuo piano di stabilire un limite di 2 turni per continuare a giocare. Se non riusciamo a migliorare significativamente la nostra posizione entro i prossimi 2 turni, dichiariamo la fine del gioco. Sembra una strategia ragionevole considerando il numero di penalità accumulate e il tempo rimanente. Propongo di procedere con il prossimo turno e vedere se possiamo fare progressi. Se non miglioriamo, potremo valutare di dichiarare la fine del gioco al turno successivo. ' 
	  } 
	   } 
	   } 
	 & & \\ 
 

    \theutterance \stepcounter{utterance}  

    & & \multicolumn{4}{p{0.6\linewidth}}{\cellcolor[rgb]{0.9,0.9,0.9}{%
	\makecell[{{p{\linewidth}}}]{% 
	  \tt {\tiny [A$\langle$GM]}  
	 Penalità: Il tuo messaggio non deve contenere nulla prima del comando!\\ \tt \\ \tt Assicurati che la tua risposta contenga solo `COMUNICA: <MESSAGE>` o `SPOSTA: <OBJECT>, (<X>, <Y>)`, e nient'altro!\\ \tt \\ \tt You have collectively accumulated 12 of 12 penalties.\\ \tt Riprova! 
	  } 
	   } 
	   } 
	 & \\ 
 

    \theutterance \stepcounter{utterance}  

    & \multicolumn{4}{p{0.6\linewidth}}{\cellcolor[rgb]{0.8,1,0.9}{%
	\makecell[{{p{\linewidth}}}]{% 
	  \tt {\tiny [A$\rangle$GM]}  
	 COMUNICA: finito? 
	  } 
	   } 
	   } 
	 & & \\ 
 

    \theutterance \stepcounter{utterance}  

    & & & \multicolumn{2}{p{0.3\linewidth}}{\cellcolor[rgb]{0.95,0.95,0.95}{%
	\makecell[{{p{\linewidth}}}]{% 
	  \tt {\tiny [GM$|$GM]}  
	 COMUNICA: finito? 
	  } 
	   } 
	   } 
	 & & \\ 
 

    \theutterance \stepcounter{utterance}  

    & & \multicolumn{4}{p{0.6\linewidth}}{\cellcolor[rgb]{0.9,0.9,0.9}{%
	\makecell[{{p{\linewidth}}}]{% 
	  \tt {\tiny [GM$\rangle$B]}  
	 Il tuo messaggio è stato inoltrato all'altro giocatore.\\ \tt \\ \tt \\ \tt You are currently playing round 19 of maximum 20.\\ \tt You have collectively accumulated 12 of 12 penalties.\\ \tt L'altro giocatore ha inviato questo messaggio:\\ \tt \\ \tt \\ \tt "finito?"\\ \tt \\ \tt \\ \tt Qual è il tuo prossimo comando? 
	  } 
	   } 
	   } 
	 & \\ 
 

    \theutterance \stepcounter{utterance}  

    & & & \multicolumn{4}{p{0.6\linewidth}}{\cellcolor[rgb]{1,0.85,0.72}{%
	\makecell[{{p{\linewidth}}}]{% 
	  \tt {\tiny [GM$\langle$B]}  
	 COMUNICA: finito! 
	  } 
	   } 
	   } 
	  \\ 
 

    \theutterance \stepcounter{utterance}  

    & & & \multicolumn{2}{p{0.3\linewidth}}{\cellcolor[rgb]{0.95,0.95,0.95}{%
	\makecell[{{p{\linewidth}}}]{% 
	  \tt {\tiny [GM$|$GM]}  
	 COMUNICA: finito! 
	  } 
	   } 
	   } 
	 & & \\ 
 

    \theutterance \stepcounter{utterance}  

    & & & \multicolumn{2}{p{0.3\linewidth}}{\cellcolor[rgb]{0.95,0.95,0.95}{%
	\makecell[{{p{\linewidth}}}]{% 
	  \tt {\tiny [GM$|$GM]}  
	 true 
	  } 
	   } 
	   } 
	 & & \\ 
 

    \theutterance \stepcounter{utterance}  

    & & & \multicolumn{2}{p{0.3\linewidth}}{\cellcolor[rgb]{0.95,0.95,0.95}{%
	\makecell[{{p{\linewidth}}}]{% 
	  \tt {\tiny [GM$|$GM]}  
	 * success: True\\ \tt * lose: False\\ \tt * aborted: False\\ \tt -------\\ \tt * Shifts: 2.00\\ \tt * Max Shifts: 8.00\\ \tt * Min Shifts: 4.00\\ \tt * End Distance Sum: 9.43\\ \tt * Init Distance Sum: 20.21\\ \tt * Expected Distance Sum: 20.95\\ \tt * Penalties: 12.00\\ \tt * Max Penalties: 12.00\\ \tt * Rounds: 19.00\\ \tt * Max Rounds: 20.00\\ \tt * Object Count: 5.00\\ \tt  
	  } 
	   } 
	   } 
	 & & \\ 
 

    \theutterance \stepcounter{utterance}  

    & & & \multicolumn{2}{p{0.3\linewidth}}{\cellcolor[rgb]{0.95,0.95,0.95}{%
	\makecell[{{p{\linewidth}}}]{% 
	  \tt {\tiny [GM$|$GM]}  
	 * Main Score: 1.68\\ \tt \\ \tt -------\\ \tt * Distance Score: 0.54\\ \tt * Coverage Score: 0.11\\ \tt * Penalty Score: 0.09\\ \tt * Alternative Penalty Score: 0.17\\ \tt * Alternative Main Score: 9.02\\ \tt \\ \tt -------\\ \tt * Shifts: 2.00\\ \tt * Max Shifts: 8.00\\ \tt * Min Shifts: 4.00\\ \tt * End Distance Sum: 9.43\\ \tt * Init Distance Sum: 20.21\\ \tt * Expected Distance Sum: 20.95\\ \tt * Penalties: 12.00\\ \tt * Max Penalties: 12.00\\ \tt * Rounds: 19.00\\ \tt * Max Rounds: 20.00\\ \tt * Object Count: 5.00\\ \tt  
	  } 
	   } 
	   } 
	 & & \\ 
 

\end{supertabular}
}

\end{document}
